\documentclass{ctexart}

% === 页面布局 ===
\usepackage[left=2.5cm, right=2.5cm, top=2.5cm, bottom=2.5cm]{geometry}

% === 数学与符号 ===
\usepackage{amsmath, amssymb}

% === 链接 ===
\usepackage[colorlinks,linkcolor=black,citecolor=black,urlcolor=black]{hyperref}

% === 定理类环境 ===
\newtheorem{theorem}{定理}[section]
\newtheorem{lemma}[theorem]{引理}
\newtheorem{corollary}[theorem]{推论}
\newtheorem{definition}[theorem]{定义}
\newtheorem{conjecture}[theorem]{猜想}
\newtheorem{axiom}[theorem]{公理}
\newtheorem{remark}[theorem]{注记}
\newtheorem{proof}[theorem]{证明}

\begin{document}
	
	\title{\bfseries 不可公度性、复布尔图灵机与实布尔图灵机：\\
		非确定性计算的代数语义框架}
	\author{郑波尽\thanks{中南民族大学计算机学院，武汉，430074，中国。}
		\and 郑景文\thanks{武汉大学网络安全学院，武汉，430072，中国。}
		\and 王伟武\thanks{国家开发银行信息科技处，北京，100033，中国。}
		\and 黄芷莹\thanks{中南民族大学计算机学院，武汉，430074，中国。}
	}
	\date{2026年7月}
	
	\maketitle
	
	\begin{abstract}
		本文构建了非确定性计算的代数语义框架，由三个逻辑上紧密关联的部分组成。第一部分严格证明了一个关于经典图灵机模型表达能力的根本性命题：不可公度性，即两个代数数在基域上线性无关的性质，不能是经典图灵机操作语义的内在属性。该证明基于语法不变性原理：图灵机的转移函数仅依赖符号标识，因此机器行为在符号置换下同构；任何操作语义的内在属性必须在此置换下保持不变。通过构造具体的符号置换和语义重赋值，不可公度性随外部解释而变，故不可能是内在属性。这一定理构成了整个维度退化理论的元理论基石，直接解释了为什么经典NP定义将非确定性的本质外包给外部量词``存在见证''和外部谓词$IsWitness(c)$——因为经典确定性图灵机的内部语义无法表达刻画非确定性分支独立性的核心概念。在此基础上，我们进一步诊断了经典NP定义的操作不自足与三重概念错误：经典NP定义隐式要求了一个在经典操作语义中不可实现的虚部到实部拷贝操作；由此必然导致第一重错误（函数指针被强制转换为数据，同时导致不可公度性的丢失）、第二重错误（对偶性被存在量词消除）和第三重错误（因不可公度性和对偶性皆已丢失，只能将非确定性的本质外包，导致语义残缺）。
		
		基于这一诊断，第二部分构建了复布尔图灵机（CBTM），将计算符号从布尔值提升为四元伽罗瓦域$\mathbb{F}_4 = \{0, 1, \alpha, \beta\}$中的代数元素。CBTM的核心洞见在于：非确定性计算对应于代数域的扩张——当读取虚部为$1$的符号时，计算必须分叉为两条路径，分别绑定于在基域$\mathbb{F}_2$上不可公度且互为对偶的生成元$\alpha$和$\beta$。通过投影算子$\operatorname{Re}$和$\operatorname{Im}$，我们分离``旧数据''与``新维度''，并引入双带视角将抽象的代数符号直观地分解为实带（确定性计算）和虚带（非确定性控制）。参数化等价定理严格证明了CBTM与经典模型的外延等价（$\mathbf{P}_{cb}=\mathbf{P}$且$\mathbf{NP}_{cb}=\mathbf{NP}$），但其操作语义具有经典模型所不具备的内涵：CBTM内建了不可公度性，显式保留了对偶性，从根本上消除了拷贝操作的必要性，以语义内建取代了语义外包。由此，CBTM框架为经典非确定性图灵机补充了长期缺失的前提——多个选项之所以构成真正的非确定性选择，是因为它们对应于域扩张中相互不可公度且互为对偶的元素——并系统性地修正了经典NP定义的操作不自足和三重概念错误。参数化猜想将$\mathbf{P}$ vs.\ $\mathbf{NP}$等价转化为虚部$b$是否可消除的代数问题。
		
		第三部分将抽象生成元$\alpha$具体化为代数数$\sqrt{2}$，引入实布尔图灵机（RBTM）和双带视角。虚带上的``1''直观地标记了``新维度''出现的位置，为后续的动态维度跟踪奠定了物理基础。我们证明了生成元无关性定理：计算能力不依赖于生成元的具体选择，无论使用$\sqrt{2}$、$\sqrt{3}$还是虚数单位$i$，相应的自动机都是同构的。这揭示了非确定性的本质在于``引入与基域不可公度的新元素''这一事实本身，而非生成元的代数身份。基于此，我们进一步揭示了伽罗瓦自同构的操作本质：它在计算树上的作用，就是将每条路径的所有非确定性选择全部翻转，即切换到其对偶路径。在虚部验证机（IVM）的动态扩张下，伽罗瓦群扩展为$(\mathbb{Z}/2\mathbb{Z})^m$，允许独立翻转单个非确定性选择点（单点敏感性），从而为精确刻画非确定性提供操作基础。我们据此提出了敏感性猜想：$\mathbf{NP}$完全语言的验证器必然具有不可消除的单点敏感性，而$\mathbf{P}$类语言可以完全不敏感。该猜想被证明等价于$\mathbf{P} \neq \mathbf{NP}$。这一等价性将P vs.\ NP问题从定量资源分析转化为定性结构判据，为后续研究提供了新的分析路径。
		
		本文所建立的框架为后续论文中虚部验证机（IVM）的动态维度追踪、本质维度$\kappa(L)$的定义、障碍无关性论证以及$\mathbf{P}\neq\mathbf{NP}$的严格分析提供了完整的模型基础。
		
		\vspace{0.5cm}
		\noindent \textbf{关键词：} 不可公度性，语法不变性，对偶性，操作不自足，复布尔图灵机，实布尔图灵机，代数语义，域扩张，非确定性计算，参数化猜想，生成元无关性，伽罗瓦自同构，敏感性猜想
	\end{abstract}
	
	\section{引言}
	
	\subsection{P vs.\ NP问题与三大障碍}
	
	$\mathbf{P}$ versus $\mathbf{NP}$问题\cite{cook1971,levin1973}是理论计算机科学乃至整个现代数学中最具影响力的未解难题之一。复杂度类$\mathbf{P}$由所有能够在多项式时间内被确定性图灵机判定的判定问题组成；复杂度类$\mathbf{NP}$由所有能够在多项式时间内被非确定性图灵机判定的判定问题组成。$\mathbf{P} \subseteq \mathbf{NP}$是平凡的，但$\mathbf{P}$是否等于$\mathbf{NP}$——即所有可以在多项式时间内被验证的问题是否也可以在多项式时间内被求解——至今未解。若$\mathbf{P} = \mathbf{NP}$，则现代密码体系将面临严峻挑战；若$\mathbf{P} \neq \mathbf{NP}$，则将确认计算世界中存在内在的复杂性层级。
	
	经过半个多世纪的努力，研究者们逐渐识别出三个经典的证明障碍。Baker、Gill和Solovay~\cite{baker1975}的相对化障碍证明，存在两个谕示世界，一个使$\mathbf{P}=\mathbf{NP}$，另一个使$\mathbf{P}\neq\mathbf{NP}$，任何对谕示``透明''的证明技术都无法同时兼容这两个世界。Razborov和Rudich~\cite{razborov1997}的自然证明障碍指出，任何同时满足``有用性''、``大尺度''和``构造性''的组合下界证明将导致伪随机数生成器的攻破，因而被认为几乎不可能。Aaronson和Wigderson~\cite{aaronson2009}的代数化障碍将相对化推广为代数化，证明任何对谕示多项式扩展保持有效的证明技术同样无法分离$\mathbf{P}$和$\mathbf{NP}$。
	
	这三大障碍共同指向一个深层问题：传统计算模型是``语义贫乏''的——它们只关心符号的语法操作，抛弃了符号背后的数学语义。语义的缺失使得任何纯语法的论证都容易被外部谕示或伪随机函数所愚弄。障碍的存在暗示，问题的难度可能不仅来自计算本身，更来自我们用以描述计算的数学语言。
	
	\subsection{经典模型的语义贫乏：从语法到语义的鸿沟}
	
	经典图灵机模型的核心特征是彻底的\textbf{语法中立性}。一台图灵机是一个七元组$M = (Q, \Sigma, \Gamma, \delta, q_0, q_{\text{accept}}, q_{\text{reject}})$，其中转移函数$\delta: Q \times \Gamma \to Q \times \Gamma \times \{L, R\}$是整个模型的灵魂。转移函数的核心特征是：其输入中的符号参数仅仅是$\Gamma$中的一个\textbf{标识}。$\delta$无法区分两个具有不同外部语义但形式上可互换的符号，它只能根据当前状态和当前符号的标识来决定下一步动作。同一个二进制串``101''可以同时被解释为数字5、布尔序列``真、假、真''或某个代数扩张中的元素$a+b\sqrt{2}$的编码。图灵机不关心这些符号的数学含义，它只是机械地根据转移函数移动符号，计算的过程是符号操作，但计算的``意义''完全存在于外部观察者的解释之中。
	
	正是这种语法中立性赋予了图灵机通用模拟的能力——它可以模拟任何算法，因为它不预设任何特定的语义。但也正是这种语法中立性，埋下了一个深层的隐患：当我们需要区分计算的``难度''时，纯粹语法层面的描述可能丢失了决定难度的关键数学信息。
	
	在经典计算理论中，$\mathbf{P}$的定义是内在的、操作性的：存在多项式时间确定性图灵机。$\mathbf{P}$的语义完全由机器的操作语义提供。而$\mathbf{NP}$的定义是外在的、外延性的：
	\[
	x \in L \iff \exists c \; [IsWitness_L(c) \land (|c| \le p(|x|) \land V(x, c) = 1)].
	\]
	其中$IsWitness_L(c)$是一个外部谓词，断言``$c$是关于语言$L$的一个见证''。在经典框架中，这个谓词的语义被外延性地还原为验证器的接受行为。但见证与输入之间的内在代数关系——即非确定性分支选项之间的独立性——从未被正面定义。
	
	经典非确定性图灵机（NTM）的转移关系允许在同一个（状态，符号）对下存在多个可能的后续配置。然而，对这种``多值性''的解释长期依赖于随机选择或全部分支并行探索的直觉隐喻。一个更深层的问题长期被忽视：在标准NTM的转移关系中，我们写下多个五元组作为``不同的选择''，但它们为何是不同的？两个选项之间的``独立性''由什么来保证？标准定义对此保持沉默——它仅以语法的形式列举选项，而将``它们构成真正的非确定性选择''这一事实视为不言自明的直觉前提。非确定性的本质被外包给了元语言的存在量词``存在见证''（$\exists c$）和外部谓词$IsWitness(c)$。
	
	更为严重的是，从本文所提出的CBTM框架出发，可以发现经典NP定义中存在更深层的缺陷：经典NP定义隐式地要求了一个在经典操作语义中\textbf{不可实现}的操作——将虚部控制信息拷贝到实部成为数据$c$，这使得经典NP定义是\textbf{操作不自足}的。由此必然导致三重概念错误：
	\begin{enumerate}
		\item \textbf{函数指针被强制转换为数据，同时导致不可公度性的丢失}：见证$c$本质上是控制流信息（函数指针地址序列），但在经典NP定义中被写入纸带，被强制当作数据来处理。更为严重的是，这一转换使得见证$c$的各个比特之间原本的不可公度性完全消失——在虚部中，不同比特对应于不同生成元的选择（如$\sqrt{p_2}$和$\sqrt{p_3}$），这些生成元在基域上线性无关，保证了选择路径的独立性；转换为数据比特$0$和$1$后，比特间不再保留任何代数独立性关系。
		\item \textbf{对偶性被存在量词消除}：在非确定性分支中，每一次选择都是在两个对偶的生成元之间做出的（$\alpha$与$\beta$，$\sqrt{p_i}$与$-\sqrt{p_i}$）。当这些选择被编码为比特串$c$并用存在量词$\exists c$包裹时，对偶性被完全抹去——那个未被选择的路径的所有痕迹都消失了。
		\item \textbf{语义外包导致概念残缺}：因为不可公度性已在第一重错误中丢失，对偶性已在第二重错误中丢失，经典模型无法正面刻画非确定性选择的独立性和结构，只能将其外包给外部量词和外部谓词。
	\end{enumerate}
	
	本文的核心诊断是：\textbf{经典确定性图灵机的内部语义无法内在地表达刻画非确定性分支独立性的核心概念——不可公度性。}不可公度性，即两个代数数在基域上线性无关的性质，是保证两个非确定性选择路径在代数上不可相互归约的数学基础。经典图灵机的操作语义是语法性的，它只能表达基域内的仿射变换，而无法``感知''这种跨越基域边界的代数关系。与此同时，经典NP定义还在外包过程中主动消除了非确定性的对偶结构，并在编码过程中丢失了不可公度性。这一双重语义缺失——不可公度性在编码中丢失，对偶性在量词中被消除——正是三大障碍能够长期存在的根本原因。
	
	\subsection{核心洞见：非确定性作为代数扩张}
	
	本文的核心洞见是：\textbf{非确定性选择，本质上就是在基域中引入一个与旧元素不可公度的新元素。}考虑有理数域$\mathbb{Q}$和它的二次扩张$\mathbb{Q}(\sqrt{2})$。$\sqrt{2}$与$\mathbb{Q}$中的数``不可公度''——它不能表示为任何有理数的有限线性组合。当我们面对一个涉及$\sqrt{2}$的代数问题时，我们无法用一条确定性的路径同时处理``包含$\sqrt{2}$''和``不包含$\sqrt{2}$''两种情况；我们必须分叉，同时探索两个世界。这正是非确定性的本质。而且，这两个分支选择不是任意的——它们互为对偶：$\sqrt{2}$与$-\sqrt{2}$满足完全相同的代数方程$x^2-2=0$，由伽罗瓦自同构精确地相互映射。
	
	将这一思想抽象化：我们用二元域$\mathbb{F}_2$（对应布尔值0/1）作为基域，它的二次扩张$\mathrm{GF}(4)$恰好引入了一个新元素$\alpha$（满足$\alpha^2 = \alpha + 1$）。$\mathrm{GF}(4)$中的每个元素都可以唯一地表示为$a + b\alpha$，其中$a,b \in \{0,1\}$。这里的$a$代表``旧数据''，$b$指示``是否涉及新维度''。当$b=1$时，符号携带着新维度的信息，计算必须分叉，而分叉的两条路径分别绑定于对偶生成元$\alpha$和$\beta = \alpha+1$。
	
	这一核心洞见将非确定性从``有多个选项''的语法约定提升为``引入不可公度元并强制分叉为对偶路径''的代数事实，为我们构建新的计算模型提供了根本动机。
	
	\subsection{本文的贡献与结构}
	
	本文从三个逻辑上紧密关联的层面展开上述诊断并构建相应的解决方案。
	
	\textbf{第一，元理论基石（第3节）。}我们基于语法不变性原理，严格证明了不可公度性不能是经典图灵机操作语义的内在属性（定理\ref{thm:main}）。该证明的核心逻辑是：图灵机的转移函数仅依赖符号标识，机器行为在符号置换下同构，因此任何操作语义的内在属性必须在符号置换下保持不变。通过构造具体的符号置换和语义重赋值，不可公度性在一种解释下存在，在另一种解释下消失，故不可能是内在属性。这一定理构成了整个维度退化理论的元理论基石，直接解释了为什么经典NP定义将非确定性的本质外包给外部量词和外部谓词，并在此过程中导致了不可公度性的丢失和对偶性的消除。
	
	\textbf{第二，复布尔图灵机（第4节）。}基于上述诊断，我们构建了复布尔图灵机（CBTM），将纸带符号提升为四元有限域$\mathbb{F}_4$中的代数元素。CBTM的核心洞见在于：非确定性计算对应于代数域的扩张。通过投影算子$\operatorname{Re}$和$\operatorname{Im}$，我们将每个符号分解为实部（旧数据）和虚部（新维度标记）。参数化等价定理严格确立了CBTM与经典模型的外延等价（$\mathbf{P}_{cb}=\mathbf{P}$，$\mathbf{NP}_{cb}=\mathbf{NP}$），同时揭示了CBTM操作语义的内涵增强——它内建了不可公度性，显式保留了对偶性，从根本上消除了拷贝操作的必要性，系统性地修正了经典NP定义的操作不自足和三重概念错误。参数化猜想将$\mathbf{P}$ vs.\ $\mathbf{NP}$等价转化为虚部$b$是否可消除的代数问题。
	
	\textbf{第三，实布尔图灵机与生成元无关性（第5节）。}我们将抽象生成元$\alpha$具体化为代数数$\sqrt{2}$，引入实布尔图灵机（RBTM）和双带视角。生成元无关性定理证明了计算能力不依赖于生成元的具体选择，揭示了非确定性的本质在于``引入与基域不可公度的新元素''这一事实本身。我们揭示了伽罗瓦自同构的操作性本质——路径切换——并据此提出敏感性猜想：P与NP的区分等价于是否存在某个非确定性选择点，使得单独翻转该点的选择能够改变计算结果。该猜想等价于$\mathbf{P}\neq\mathbf{NP}$。
	
	\textbf{第四，方法论原则（第6节）。}我们提出``数计一体''的方法论原则——计算理论应当仿照数学，要求任意算法都应当具有明确的语义，并且语义与语法保持一致。CBTM/RBTM框架正是这一原则的严格实现，可以精确地理解为对经典NP定义的操作不自足和三重概念错误的系统性修正。
	
	第7节总结本文贡献并展望后续工作。
	
	\section{预备知识}
	
	\subsection{经典图灵机与非确定性}
	
	\begin{definition}[经典确定性图灵机]\label{def:dtm}
		一台经典（确定性）图灵机是一个七元组$M = (Q, \Sigma, \Gamma, \delta, q_0, q_{\text{accept}}, q_{\text{reject}})$，其中$Q$是有限状态集，$\Sigma \subseteq \Gamma$是输入字母表，$\Gamma$是有限纸带字母表，$\delta: Q \times \Gamma \to Q \times \Gamma \times \{L, R\}$是（部分）转移函数，$q_0 \in Q$是起始状态，$q_{\text{accept}}, q_{\text{reject}} \in Q$分别是接受和拒绝状态。字母表$\Gamma$包含一个特殊的空白符$B \in \Gamma \setminus \Sigma$。初始格局中，输入字符串写在纸带上，其余所有单元为空白符，带头位于最左非空符号处，状态为$q_0$。
	\end{definition}
	
	转移函数$\delta$的核心特征是：其输入中的符号参数仅仅是$\Gamma$中的一个\textbf{标识}。$\delta$无法区分两个具有不同外部语义但形式上可互换的符号，它只能根据当前状态和当前符号的标识来决定下一步动作。这一语法中立性是图灵机通用性的基石，也是本文所揭示的经典模型表达局限的根源。
	
	\begin{definition}[经典非确定性图灵机]\label{def:ntm}
		一台经典非确定性图灵机是一个七元组$N = (Q, \Sigma, \Gamma, \Delta, q_0, q_{\text{accept}}, q_{\text{reject}})$，其中$Q, \Sigma, \Gamma, q_0, q_{\text{accept}}, q_{\text{reject}}$的含义与DTM相同，而转移关系$\Delta \subseteq Q \times \Gamma \times Q \times \Gamma \times \{L, R\}$允许在同一个（状态，符号）对下存在多个五元组。$N$的计算是一棵配置树。如果至少有一个分支到达$q_{\text{accept}}$，则接受输入；如果所有分支都到达$q_{\text{reject}}$，则拒绝输入。
	\end{definition}
	
	\begin{definition}[复杂度类]\label{def:complexity}
		$\mathbf{P}$是所有可由多项式时间确定性图灵机判定的语言类。$\mathbf{NP}$是所有可由多项式时间非确定性图灵机判定的语言类。等价地，$L \in \mathbf{NP}$当且仅当存在多项式$p$和多项式时间确定性图灵机$V$（验证器），使得$x \in L \iff \exists c \; [IsWitness_L(c) \land |c| \le p(|x|) \land V(x,c) = 1]$。
	\end{definition}
	
	\subsection{代数背景：域、扩张与伽罗瓦群}
	
	\begin{definition}[域扩张]
		设$K \subseteq L$为两个域。若$L$作为$K$上的向量空间是有限维的，则称$L/K$为有限扩张，维数记为$[L:K]$。
	\end{definition}
	
	\begin{definition}[不可公度性]\label{def:algebraic-incomm}
		设$V$为域$K$上的向量空间。两个非零向量$u, v \in V$称为在$K$上不可公度，如果它们在$K$上线性无关。在本文中，我们主要关注代数数在有理数域$\mathbb{Q}$上的不可公度性：两个代数数$\gamma, \delta$在$\mathbb{Q}$上不可公度，当且仅当不存在不全为零的有理数$c_1, c_2$使得$c_1\gamma + c_2\delta = 0$。这等价于$\gamma$与$\delta$皆非零，且其比值$\gamma/\delta$为无理数。
	\end{definition}
	
	对于两个元素，``不可公度''与``线性无关''是等价的。本文定理直接处理两个元素的情形，因此使用``不可公度性''一词。对于多个元素，``线性无关性''是``不可公度性''的自然推广。在后续论文关于虚部验证机（IVM）和维度下界的讨论中，我们将看到线性无关性在多步分支场景下的核心作用。
	
	这一定义确保了零元素不参与不可公度性的判定：包含零的集合天然是线性相关的（可公度）。
	
	\begin{definition}[伽罗瓦扩张与伽罗瓦群]\label{def:galois}
		有限扩张$L/K$称为伽罗瓦扩张，如果$L$是$K$上某个可分多项式的分裂域。伽罗瓦群$\operatorname{Gal}(L/K)$是保持$K$逐点不动的$L$的自同构群。
	\end{definition}
	
	在本文中，我们主要使用两个具体的伽罗瓦扩张。第一个是有限域$\mathbb{F}_4$对$\mathbb{F}_2$的扩张：多项式$x^2 + x + 1$在$\mathbb{F}_2$上不可约，其分裂域$\mathbb{F}_4 = \mathbb{F}_2(\alpha) = \{0, 1, \alpha, \beta = \alpha+1\}$满足$[\mathbb{F}_4:\mathbb{F}_2] = 2$，$\operatorname{Gal}(\mathbb{F}_4/\mathbb{F}_2) \cong \mathbb{Z}/2\mathbb{Z}$，唯一的非平凡自同构为$\sigma: \alpha \leftrightarrow \beta$。第二个是无限域$\mathbb{Q}(\sqrt{2})$对$\mathbb{Q}$的扩张，其伽罗瓦群同为$\mathbb{Z}/2\mathbb{Z}$，非平凡自同构为$\sqrt{2} \leftrightarrow -\sqrt{2}$。
	
	在后续章节中，我们将看到这两个扩张如何在非确定性计算的形式化中扮演核心角色：$\mathbb{F}_4/\mathbb{F}_2$为CBTM提供了最简单的离散代数结构，而$\mathbb{Q}(\sqrt{2})/\mathbb{Q}$则为RBTM的双带视角和生成元无关性提供了直观的几何基础。两者虽然在代数结构上并不同构（特征不同），但在布尔投影层所诱导的自动机却是同构的——这正是生成元无关性定理的核心贡献之一。
	
	
	\section{不可公度性不是操作语义的内在属性}
	
	本节给出本文的第一个核心定理及其完整证明。该定理构成了整个维度退化理论的元理论基石：经典图灵机无法内在地表达不可公度性。该证明不依赖于CBTM/IVM框架的任何结论，是一个独立的数学结果。我们从经典图灵机最基本的语法性质出发，逐步建立严格的形式化论证。
	
	\subsection{语法不变性：经典操作语义的核心}
	
	经典图灵机的转移函数$\delta$仅依赖符号的标识（形状），而非符号可能携带的任何外部语义。这一事实可以形式化为以下引理。我们首先回顾经典图灵机的标准定义（见定义\ref{def:dtm}），然后在此基础上展开论证。
	
	\begin{lemma}[语法不变性引理]\label{lem:syntactic-invariance}
		设$M = (Q, \Gamma, \delta, q_0, q_{\text{accept}}, q_{\text{reject}})$为一台经典图灵机，其中$\Gamma$包含空白符$B$。设$\pi: \Gamma \to \Gamma$为字母表上的置换（双射），满足$\pi(B) = B$（即置换保持空白符不变）。构造新机器$M_\pi = (Q, \Gamma, \delta_\pi, q_0, q_{\text{accept}}, q_{\text{reject}})$，其中对于所有$q \in Q, c \in \Gamma$：
		\[
		\delta_\pi(q, \pi(c)) = (q', \pi(c'), D) \iff \delta(q, c) = (q', c', D).
		\]
		则$M$与$M_\pi$结构同构：若$M$在输入$w = w_1 w_2 \ldots w_n$上经历格局序列$C_0, C_1, C_2, \ldots$，则$M_\pi$在输入$\pi(w) = \pi(w_1) \pi(w_2) \ldots \pi(w_n)$上经历格局序列$\pi(C_0), \pi(C_1), \pi(C_2), \ldots$，其中$\pi(C)$表示将格局$C$纸带上所有符号$c$替换为$\pi(c)$所得到的格局。
	\end{lemma}
	\begin{proof}
		对计算步数进行归纳。条件$\pi(B)=B$保证$M_\pi$的初始格局是合法的（非输入部分仍为空白符）。归纳基础：初始格局$C_0$与$\pi(C_0)$对应。归纳步骤：若第$t$步$M$在格局$C_t$处读取符号$c$并转移到$C_{t+1}$，则$M_\pi$在格局$\pi(C_t)$处读取符号$\pi(c)$，根据$\delta_\pi$的定义，执行相同的动作并转移到$\pi(C_{t+1})$。
	\end{proof}
	
	语法不变性引理揭示了一个基本事实：只要不置换空白符，图灵机的行为在数据符号重标记下是完全对称的。将所有的``0''换成``1''、所有的``1''换成``0''，机器的操作轨迹不变，仅仅是符号的标签变了。这种对称性并非图灵机的缺陷——恰恰相反，它是图灵机通用性的基石，因为它意味着机器不预设任何特定的语义解释。但这也意味着，任何机器能够``内在地''区分或响应的属性，必须在符号重标记下保持不变。这一观察构成了我们定义``内在属性''并证明主要定理的逻辑基础。
	
	\subsection{内在属性与置换不变性}
	
	\begin{definition}[内在属性]\label{def:intrinsic}
		设$\mathfrak{P}$是机器$M$在输入$w$上的一个属性。称$\mathfrak{P}$是图灵机操作语义的\textbf{内在属性}，如果它完全由机器的转移规则和符号标识决定，而不依赖于外部观察者对符号的语义赋值。形式化地，$\mathfrak{P}$是内在属性当且仅当对于字母表上任意保持空白符不变的置换$\pi: \Gamma \to \Gamma$，$M$在输入$w$上具有$\mathfrak{P}$当且仅当$M_\pi$在输入$\pi(w)$上具有$\mathfrak{P}$。
	\end{definition}
	
	这一形式化刻画基于经典图灵机操作语义的核心特征：转移函数$\delta$仅依赖当前状态和读写头下符号的标识，而不依赖任何外部语义。因此，机器的全部行为完全由状态集$Q$、字母表$\Gamma$、以及基于这些集合定义的转移规则$\delta$所决定。在此设定下，任何能够被机器``内在地''区分或响应的属性，必须在字母表$\Gamma$的重命名（即置换）下保持不变，因为重命名仅仅改变了符号的标签，而没有改变状态结构和转移规则。内在属性的置换不变性因此穷尽了基于语法规则本身所能区分的所有属性，是经典图灵机语法本质的直接体现。
	
	例如，``机器在输入$w$上停机''是一个内在属性——无论我们将符号解释成什么数学对象，停机的判断只取决于操作规则本身。以下推论是语法不变性引理的直接结果，也是我们证明核心定理的关键工具。
	
	\begin{corollary}[内在属性的置换不变性]\label{cor:intrinsic}
		若属性$\mathfrak{P}$是机器$M$操作语义的内在属性，则对于任意保持空白符不变的置换$\pi: \Gamma \to \Gamma$，$M$在输入$w$上具有$\mathfrak{P}$当且仅当$M_\pi$在输入$\pi(w)$上具有$\mathfrak{P}$。
	\end{corollary}
	
	其逆否命题是本文证明的核心逻辑：\textbf{如果一个属性在某个置换下发生了改变，那么该属性必然依赖于外部语义赋值，不可能是操作语义的内在属性。}这一简单而强大的逻辑使我们能够通过构造一个精巧的符号置换来证明不可公度性不是内在属性。
	
	\subsection{两种意义上的``知道''：算法可判定性与操作语义内蕴性}
	
	在继续论证之前，我们需要区分两种意义上的``知道''一个数学属性。这一区分对于精确理解元定理的结论及其对P vs.\ NP问题的意义至关重要。
	
	在\textbf{弱意义}（算法可判定性）上，我们说一台机器``知道''一个数学属性，如果它在给定适当编码的输入后，能够输出关于该属性的正确判定。例如，经典图灵机可以在适当的编码下判定两个代数数是否在$\mathbb{Q}$上线性无关——通过高斯消元等多项式时间算法。不可公度性本身在$\mathbf{P}$内是可计算的，这是一个重要的数学事实。然而，这种``知道''依赖于外部编码和语义赋值：编码函数将数学对象映射为符号串，机器操作这些符号串，结果再由外部观察者解码回数学命题。机器本身并不``理解''它所操作的对象具有不可公度性——它只是在执行语法规则。
	
	在\textbf{强意义}（操作语义的内蕴性）上，我们说一台机器``知道''一个数学属性，如果该属性是机器操作语义的内在属性（定义\ref{def:intrinsic}）。本文的核心论断是：\textbf{经典图灵机可以在弱意义上判定不可公度性，但无法在强意义上内在地锚定它。}元定理并不否认经典图灵机可以通过适当的编码来判定两个代数数是否可公度。它证明的是一个更强的结论：图灵机的操作语义本身无法内在地``感知''它所处理的符号是否具有不可公度性——机器的操作行为不能唯一地确定符号的外部解释必须包含不可公度元。
	
	这一区分的理论分量在于：P vs.\ NP问题的困难不在于``不可公度性是否可判定''，而在于``经典模型的操作语义是否能够正面定义非确定性选择的独立性''。后者恰恰需要不可公度性作为操作语义的内在属性——而非仅仅作为外部编码的语义赋值。我们将在第3.6节详细展开这一诊断。
	
	\subsection{编码不可公度性}
	
	为了严格地讨论``机器是否能处理不可公度元''，我们必须显式地引入外部语义。这通过一个编码函数来完成。在本理论体系中，不可公度性是对于非确定性分支独立性最自然的候选刻画：如果两个分支选项对应于在基域上线性无关的生成元，那么它们就无法通过基域内的任何运算相互转化，从而构成``真正的非确定性选择''。本节定理的核心任务是证明：即使是这种最弱意义上（仅取决于输入和编码，与机器行为无关）的不可公度性，经典图灵机也无法内在地锚定。
	
	\begin{definition}[代数不可公度性]\label{def:algebraic-incomm}
		两个代数数$\gamma, \delta$在有理数域$\mathbb{Q}$上\textbf{代数不可公度}，当且仅当它们在$\mathbb{Q}$上线性无关，即不存在不全为零的有理数$c_1, c_2$使得$c_1\gamma + c_2\delta = 0$。这等价于要求$\gamma$与$\delta$皆非零，且其比值$\gamma/\delta$为无理数。
	\end{definition}
	
	对于两个元素，``不可公度''与``线性无关''是等价的。本文定理直接处理两个元素的情形，因此使用``不可公度性''一词。对于多个元素，``线性无关性''是``不可公度性''的自然推广：多个元素线性无关蕴含了两两不可公度，但反之不成立。本文定理证明经典模型无法内在地表达不可公度性，因而也无法内在地表达线性无关性（后者是前者的多元推广）。关于线性无关性在多步分支场景下的核心作用，详见关于虚部验证机（IVM）和本质维度的后续工作。
	
	这一定义确保了零元素不参与不可公度性的判定：包含零的集合天然是线性相关的（可公度）。
	
	\begin{definition}[代数编码与编码不可公度性]\label{def:encoding}
		设$\varepsilon_0: \Gamma \setminus \{B\} \to \mathbb{A}$是一个为每个数据符号指定一个代数数的函数。将其扩展为$\varepsilon: \Gamma^* \to \mathbb{A}$，使得对于输入$w \in \Gamma^*$，$\varepsilon(w) = \{\varepsilon_0(c) : c \text{ 是 } w \text{ 中出现的符号}\}$。称机器$M$在输入$w \in \Gamma^*$上\textbf{关联$\varepsilon$不可公度元}，如果$\varepsilon(w)$中包含两个非零代数数$\gamma, \delta$，它们在有理数域$\mathbb{Q}$上代数不可公度。
	\end{definition}
	
	注意，此处$\varepsilon(w)$被定义为符号值构成的集合而非多重集，因为线性相关性仅取决于元素的存在性，与重复次数无关。更重要的是，这一定义将``关联不可公度元''界定为一个\textbf{仅依赖于输入字符串$w$和编码$\varepsilon$、而与机器$M$的计算行为无关}的外部属性。这意味着我们考察的是一个比``机器在计算中利用了不可公度元''更弱的条件，即仅仅是``输入中包含了不可公度元''。如果经典模型连这种最弱意义上的不可公度性都无法内在地表达，那么任何更强的、涉及机器行为的定义就更加不可能。这一设计使定理的结论更强而非更弱。
	
	\subsection{主要元定理及其完整证明}
	
	为聚焦于操作语义的本质而非编码细节，下文考虑具备足够表达能力的字母表。在P vs.\ NP等标准复杂性分析语境下，通常默认允许多带图灵机或多符号字母表，因为它们不影响多项式时间等价类的定义。因此，本文的结论适用于这一标准分析语境下的图灵机变体。
	
	\begin{theorem}[不可公度性不能被内在化]\label{thm:main}
		设经典图灵机的字母表$\Gamma$包含空白符$B$及至少四个不同于$B$的数据符号。则不存在一个操作语义的内在属性$\mathfrak{P}$，使得对于所有满足以下条件的编码函数$\varepsilon$——即$\varepsilon$将每个数据符号独立地映射到某个代数数，且将任意两个指定的数据符号映射到非零有理数$a$和$b$（其中$a/b \in \mathbb{Q} \setminus \{0\}$）——和所有输入$w \in \Gamma^*$（其中$w$不包含空白符$B$），$\mathfrak{P}(M, w)$成立当且仅当$M$在$w$上关联$\varepsilon$不可公度元。等价地，经典图灵机无法通过其操作语义唯一地``锚定''不可公度性。
	\end{theorem}
	
	\noindent\textbf{注：}即使限定于三符号$\{0,1,B\}$模型，通过标准配对函数可将$k$个符号压缩为$\lceil\log_2 k\rceil$比特，所得多带模拟在多项式时间内等价，故定理结论适用于$\mathbf{P}/\mathbf{NP}$的标准分析语境。核心要点在于：问题的根源在于转移函数的语法性，而非字母表的具体大小。
	
	\begin{proof}
		设$\Gamma$的数据符号包含$\{0, 1, \alpha, \beta\}$，且它们均不同于空白符$B$。作为构造的一部分，我们选取数据符号$0$和$1$作为定理条件中所指的``两个指定符号''。我们采用反证法。
		
		\textbf{假设}：存在一个内在属性$\mathfrak{P}$，它能够精确地捕获``关联不可公度元''这一概念。即，对于任何满足条件的编码$\varepsilon$和任何由数据符号组成的输入$w$，$\mathfrak{P}(M, w)$为真当且仅当$M$在$w$上关联$\varepsilon$不可公度元。
		
		\textbf{步骤1：固定一个满足条件的编码$\varepsilon$。}
		定义编码$\varepsilon_0: \Gamma \setminus \{B\} \to \mathbb{A}$，它将数据符号映射为：
		\[
		0 \mapsto 1,\quad 1 \mapsto 2,\quad \alpha \mapsto \sqrt{2},\quad \beta \mapsto \sqrt{3},
		\]
		并扩展为$\varepsilon$。\footnote{这里有意将符号``0''和``1''映射到与自身标签不同的非零有理数（1和2），以强调编码函数仅依赖符号的标识（标签）而非符号名称的数值暗示。经典图灵机的操作语义只认识符号的标签，而不认识其被赋予的数值，这正是定理所要揭示的局限。}此编码满足定理条件：指定的两个符号被映射到非零有理数，且比值为有理数。
		
		\textbf{步骤2：构造具体输入并应用假设。}
		考虑任意一台图灵机$M$在输入$w = \alpha\beta$上的行为。根据$\varepsilon$，$\varepsilon(\alpha\beta) = \{\sqrt{2}, \sqrt{3}\}$，两者皆非零且在$\mathbb{Q}$上线性无关（不可公度）。因此，$M$在$w$上关联$\varepsilon$不可公度元。由假设，内在属性$\mathfrak{P}(M, \alpha\beta)$必须为真。
		
		\textbf{步骤3：构造符号置换$\pi$和置换后的机器$M_\pi$。}
		定义置换$\pi: \Gamma \to \Gamma$，保持空白符不变（$\pi(B)=B$），并在数据符号上的作用为：
		\[
		\pi(0) = \alpha,\quad \pi(1) = \beta,\quad \pi(\alpha) = 0,\quad \pi(\beta) = 1.
		\]
		这是一个保持空白符不变的双射。根据引理\ref{lem:syntactic-invariance}，构造$M_\pi$，它与$M$结构同构。
		
		\textbf{步骤4：应用内在属性的置换不变性。}
		由于$\mathfrak{P}$被假设为内在属性，根据推论\ref{cor:intrinsic}，它在符号置换下必须保持不变。因此，既然$\mathfrak{P}(M, \alpha\beta)$为真，那么$\mathfrak{P}(M_\pi, \pi(\alpha\beta)) = \mathfrak{P}(M_\pi, 01)$也必须为真。
		
		\textbf{步骤5：在同一编码$\varepsilon$下检验矛盾。}
		对于$M_\pi$的输入$01$，我们仍使用步骤1中固定的编码$\varepsilon$。根据$\varepsilon$，$\varepsilon(0 1) = \{1, 2\}$，两者皆为非零有理数，它们在$\mathbb{Q}$上线性相关（可公度），且不包含任何非零的不可公度元对。因此，$M_\pi$在输入$01$上并不关联$\varepsilon$不可公度元。
		
		\textbf{步骤6：导出矛盾。}
		根据步骤4，$\mathfrak{P}(M_\pi, 01)$为真。但根据步骤5和反证假设，$\mathfrak{P}(M_\pi, 01)$应为假（因为它不关联不可公度元）。这是一个矛盾。
		
		这个矛盾表明，我们最初的假设——存在一个内在属性$\mathfrak{P}$能够独立于编码地捕获关联不可公度元——是错误的。需要特别指出的是，本定理中的``关联不可公度元''被定义为纯粹的外部属性（仅取决于输入和编码），这一设计恰恰强化了结论：如果经典模型连这种最弱意义上的不可公度性都无法内在地锚定，那么任何更强的、涉及机器行为的定义就更加不可能。我们的证明用形式化手段锚定了这一直觉：任何内在属性都无法在符号置换下保持与这种外部属性的等价对应。
	\end{proof}
	
	该证明的美感在于其极简性。它不依赖于任何关于图灵机计算能力的假设，不涉及多项式时间、非确定性或任何复杂性类的概念。它仅使用了图灵机最基本的语法性质——转移函数仅依赖符号标识——以及一个精巧的符号置换构造。正是这种极简性，使得该定理具有坚实的基础，并构成整个维度退化理论的元理论基石。
	
	\subsection{元定理的意义：语义空洞、操作不自足与概念错误的揭示}
	
	元定理的直接推论是：\textbf{经典图灵机的操作语义本身永远无法``感知''不可公度性。}这不是在``算法可判定性''的弱意义上（经典图灵机可以在适当编码下判定不可公度性），而是在``操作语义内蕴性''的强意义上（它不能内在地锚定不可公度性）。无论编码多么精巧，经典图灵机的操作语义本身永远无法区分``这个符号编码了$\sqrt{2}$''和``这个符号编码了$1$''——因为符号可以被置换，而机器行为在置换下保持不变。
	
	这一定理直接解释了经典NP定义为何将非确定性的本质外包给外部量词``存在见证''（$\exists c$）和外部谓词$IsWitness(c)$。在经典框架中，$IsWitness_L(c)$的语义被外延性地还原为``存在验证器$V$使得$V(x,c)=1$''。但``是见证''这一内在属性——即$c$与$x$之间的内在代数关系，为什么$c$构成了$x$的一个解——从未被正面定义。外部谓词$IsWitness(c)$的内涵需要不可公度性来填充：见证的每一个比特代表一次独立的函数指针选择，独立性由不可公度性保证。然而，经典模型无法在强意义上表达不可公度性，因此$IsWitness_L(c)$在经典框架内是一个\textbf{语义空洞}——它指向了一个经典语言无法表达的核心概念。
	
	更进一步的诊断表明，经典NP定义还存在更深层的缺陷。该定义隐式地要求了一个在经典操作语义中不可实现的操作——将虚部控制信息拷贝到实部成为数据$c$，这使得经典NP定义是操作不自足的。由此必然导致三重概念错误：第一重错误（函数指针被强制转换为数据）同时导致不可公度性的丢失——见证$c$的各个比特在虚部中对应于不同的不可公度生成元选择，转换为数据比特后，比特间不再保留任何代数独立性关系；第二重错误导致对偶性的丢失——$+\sqrt{p_i}$与$-\sqrt{p_i}$的伽罗瓦对偶关系在比特$0$和$1$中没有对应物；因不可公度性和对偶性皆已丢失，经典模型只能将非确定性的本质外包，导致语义残缺（第三重错误）。经典框架中的``非确定性''因此是一个双重残缺的概念——既在NTM层面缺失不可公度性的正面定义，又在NP层面因编码必然丢失不可公度性和对偶性。
	
	这一语义缺失、操作不自足和概念扭曲并非疏忽，而是经典模型语言局限的必然表现。经典NP定义将非确定性的本质外包给外部量词和外部谓词，正是因为经典模型缺乏表达不可公度性所需的概念资源。三大障碍——相对化、自然证明和代数化——正是这一语义空洞、操作不自足和概念扭曲在不同技术层面的表现：相对化障碍的根源在于谕示可以替机器完成“猜测虚部控制信息并拷贝到实部”这一不可能操作，从CBTM架构来看，虚部与实部分属控制层与数据层，经典机器无法感知虚部，更无法将虚部信息跨层次拷贝到实部，但谕示可以免费提供这一拷贝服务，由此制造了$\mathbf{P}^A = \mathbf{NP}^A$的谕示世界；自然证明障碍利用的是自然性质基于组合属性，而独立性的本质是代数概念，超出了组合范畴；代数化障碍利用的是代数化谕示只能回答多项式等式问题，无法感知符号的正负号——而正负号正是对偶性在操作层面的直接体现。
	
	\subsection{代数封闭性视角的强化论证}
	
	语法不变性（置换论证）从语法层面证明了不可公度性不是内在属性。代数封闭性从另一个互补的角度强化了这一结论，使得论证更加完整。
	
	经典图灵机的操作语义，其代数表达力被限制在编码所选择的基域之内。以有理数域$\mathbb{Q}$为例：它对加法、减法、乘法、除法（非零）封闭。从有理数出发，通过有限步有理运算，永远无法``跳''到无理代数数。因此，如果编码函数将符号映射到有理数，机器通过转移函数能够操作和生成的任何信息，都永远落在$\mathbb{Q}$内。经典图灵机本质上在有理数域上进行仿射计算——每一步操作可以抽象为$X \leftarrow X + C$，其中$X$是当前状态和纸带内容的代数表示，$C$是由转移函数决定的常数。
	
	而不可公度性（如$\sqrt{2}$和$\sqrt{3}$的关系）是一种跨越基域边界的属性，它只有在扩张域$\mathbb{Q}(\sqrt{2}, \sqrt{3})$中才有定义。经典图灵机的操作语义在原则上就无法触及它。这不是``模拟协议是否聪明''的问题，而是计算模型的代数表达力边界使然。
	
	置换论证与代数封闭性论证从不同角度指向同一结论：不可公度性不可能是操作语义的内在属性。前者通过构造具体的符号置换展示不可公度性随外部解释而变，后者通过域封闭性论证揭示不可公度性是跨越基域边界的属性。两者共同为接下来的模型构建提供了坚实的理论基础：既然经典模型无法内在地表达不可公度性，我们就必须将其内建到新模型的操作语义之中。
	
	\section{复布尔图灵机：代数语义的计算模型}
	
	基于第3节的诊断——经典图灵机无法内在地表达不可公度性——本节构建复布尔图灵机（CBTM），一个将不可公度性内建为操作语义的内在属性的计算模型。CBTM的核心洞见在于：非确定性计算对应于代数域的扩张。我们将从最简单的域扩张出发，逐步建立完整的形式化体系。同时，CBTM的设计可以精确地理解为对经典NP定义的操作不自足和三重概念错误的系统性修正：从根本上消除拷贝操作的必要性，让函数指针常驻虚部（保留不可公度性），让对偶性显式保留，让语义内建而非外包。
	
	\subsection{$\mathbb{F}_4$的代数结构与投影算子}
	
	我们需要一个最简单的域扩张：从二元域$\mathbb{F}_2 = \{0,1\}$扩张到四元域$\mathbb{F}_4$。多项式$x^2 + x + 1$在$\mathbb{F}_2$上不可约。向$\mathbb{F}_2$添加该多项式的一个根$\alpha$，得到二次扩张：
	\[
	\mathbb{F}_4 = \mathbb{F}_2(\alpha) = \{0, 1, \alpha, \alpha+1\}.
	\]
	令$\beta = \alpha + 1$，则$\mathbb{F}_4 = \{0, 1, \alpha, \beta\}$。扩张次数为$[\mathbb{F}_4 : \mathbb{F}_2] = 2$，故$\mathbb{F}_4$作为$\mathbb{F}_2$上的向量空间维数为$2$，其一组基为$\{1, \alpha\}$。$\mathbb{F}_4$的特征为$2$。
	
	$\mathbb{F}_4$中元素的代数关系由$\alpha^2 + \alpha + 1 = 0$（即$\alpha^2 = \alpha + 1 = \beta$）完全确定。由此可导出\cite{lidl1997}：
	\begin{equation}\label{eq:relations}
		\alpha + \beta = 1, \quad \alpha \cdot \beta = 1, \quad \alpha^2 = \beta, \quad \beta^2 = \alpha, \quad \alpha^3 = \beta^3 = 1.
	\end{equation}
	$\alpha$与$\beta$的互逆性（$\alpha \cdot \beta = 1$）与互补性（$\alpha + \beta = 1$）精确形式化了非确定性路径之间``互为他者且不可相互归约''的对偶关系。
	
	\begin{definition}[投影算子]\label{def:projection}
		定义映射$\operatorname{Re}: \mathbb{F}_4 \to \{0,1\}$和$\operatorname{Im}: \mathbb{F}_4 \to \{0,1\}$如下：
		\[
		\begin{aligned}
			\operatorname{Re}(0) &= 0, \quad \operatorname{Im}(0) = 0; \\
			\operatorname{Re}(1) &= 1, \quad \operatorname{Im}(1) = 0; \\
			\operatorname{Re}(\alpha) &= 0, \quad \operatorname{Im}(\alpha) = 1; \\
			\operatorname{Re}(\beta) &= 1, \quad \operatorname{Im}(\beta) = 1.
		\end{aligned}
		\]
		$\operatorname{Re}(s)$称为符号$s$的\textbf{实部}，$\operatorname{Im}(s)$称为符号$s$的\textbf{虚部}。
	\end{definition}
	
	代数语义上，每个符号$s$可写为$s = \operatorname{Re}(s) + \operatorname{Im}(s) \cdot \alpha$。$\operatorname{Im}(s)=1$意味着``这个符号涉及新维度$\alpha$''，它将触发非确定性行为。这一设计将传统布尔值（实部）与``是否涉及域扩张''的标记（虚部）统一在同一个代数对象中，恢复了代码与数据的严格分离：实部是数据层，虚部是控制层。
	
	\subsection{不可公度性与伽罗瓦群}
	
	在$\mathbb{F}_4$中，取$K = \mathbb{F}_2$，$V = \mathbb{F}_4$（作为$\mathbb{F}_2$上的向量空间），则$\alpha$和$\beta$在$\mathbb{F}_2$上不可公度。不可公度性保证了两个生成元在基域内不存在任何线性关系，从而标记它们的计算路径不可通过基域内的操作相互转化。这一代数事实是CBTM刻画非确定性独立性的数学基础。
	
	扩张$\mathbb{F}_4/\mathbb{F}_2$是伽罗瓦扩张。其伽罗瓦群$\operatorname{Gal}(\mathbb{F}_4/\mathbb{F}_2)$同构于$\mathbb{Z}/2\mathbb{Z}$。唯一的非平凡自同构$\sigma: \mathbb{F}_4 \to \mathbb{F}_4$定义为：
	\begin{equation}\label{eq:sigma}
		\sigma(0) = 0, \quad \sigma(1) = 1, \quad \sigma(\alpha) = \beta, \quad \sigma(\beta) = \alpha.
	\end{equation}
	$\sigma$保持基域$\mathbb{F}_2$逐点不动，交换两个不可公度生成元。这一对偶对称性将在后续揭示伽罗瓦自同构的操作本质——路径切换——中扮演核心角色。值得注意的是，这一对偶性在经典NP定义中被存在量词消除，而在CBTM中被显式保留。
	
	\subsection{形式化定义与核心公理}
	
	\begin{definition}[复布尔图灵机]\label{def:cbtm}
		一台\textbf{复布尔图灵机}（Complex Boolean Turing Machine, CBTM）是一个七元组
		\[
		M = (Q, \Sigma, \Gamma, \Delta, q_0, q_{\text{accept}}, q_{\text{reject}}),
		\]
		其中：
		\begin{itemize}
			\item $Q$是有限状态集；
			\item $\Sigma \subseteq \mathbb{F}_4$是输入字母表；
			\item $\Gamma \subseteq \mathbb{F}_4$是纸带字母表，满足$\Sigma \subseteq \Gamma$，且存在空白符号$\sqcup \in \Gamma \setminus \Sigma$（通常取$\sqcup = 0$）；
			\item $\Delta \subseteq Q \times \Gamma \times Q \times \Gamma \times \{L,R\}$是转移关系；
			\item $q_0 \in Q$是起始状态；
			\item $q_{\text{accept}}, q_{\text{reject}} \in Q$分别是接受和拒绝状态。
		\end{itemize}
	\end{definition}
	
	机器在单条双向无限纸带上运行。初始格局为：输入字符串$w \in \Sigma^*$写在纸带上，其余单元格为空白符，带头位于最左非空符号处，状态为$q_0$。在每一步，$M$读取带头下方当前符号$\tau \in \Gamma$，根据转移关系$\Delta$决定后继动作。
	
	CBTM的行为由两条核心公理支配，它们共同构成了CBTM与经典图灵机的本质区别，并系统性地修正了经典NP定义的操作不自足和概念错误。
	
	\begin{axiom}[分支触发公理]\label{ax:branching}
		令$b = \operatorname{Im}(\tau)$。转移关系$\Delta$在$(q, \tau)$上的匹配五元组数量由$b$决定：
		\begin{enumerate}
			\item 若$b = 0$，则$\Delta$中恰好存在一个匹配的五元组。此时称机器处于\textbf{确定性模式}。符号完全在基域$\mathbb{F}_2$内，计算保持在基域中，无需代数扩张。
			\item 若$b = 1$，则$\Delta$中恰好存在两个匹配的五元组。机器产生两条后继配置，分别应用其中一个五元组。此时称机器处于\textbf{非确定性模式}。从代数角度看，两条分支分别对应于域元素$\alpha$和$\beta$所标记的不可公度路径，这两个生成元在基域上不可公度且互为对偶。\textbf{分支直接由虚部标记触发，无需将控制信息拷贝到实部，从根本上消除了经典NP定义中隐含的不可实现操作。}
		\end{enumerate}
	\end{axiom}
	
	分支触发公理将非确定性的触发机制从``转移关系中有多个五元组''的语法约定提升为``符号虚部为$1$时强制分叉为对偶路径''的代数事实。分支的二元性恰好对应$\mathbb{F}_4/\mathbb{F}_2$扩张度为$2$的代数结构，而对偶性——在经典NP定义中被存在量词消除的结构——在此被显式保留为操作语义的内在属性。
	
	\begin{axiom}[投影约束公理]\label{ax:projection}
		对于从初始格局可达的任何格局，纸带上每个单元格的内容满足$\operatorname{Re}(\tau) \in \{0,1\}$且$\operatorname{Im}(\tau) \in \{0,1\}$。写入符号的虚部由当前状态和读取符号的虚部通过布尔函数决定。特别地，若$\operatorname{Im}(\tau) = 0$，则$\operatorname{Im}(\tau') = 0$。
		
		即对CBTM的每个转移$(q',\gamma,d) \in \Delta$，写入符号$\gamma$必须满足：
		\[
		\gamma = f_{\mathfrak{Re}}(q,\mathfrak{Re}(\tau)) \oplus [\alpha \otimes f_{\mathfrak{Im}}(q,\mathfrak{Im}(\tau))],
		\]
		且若$\mathfrak{Im}(\tau)=0$，则$\mathfrak{Im}(\gamma)=0$。这里$f_{\mathfrak{Re}},f_{\mathfrak{Im}}$是取值于$\{0,1\}$的布尔函数。
	\end{axiom}
	
	投影约束公理强制任何写入纸带的新符号，其\textbf{虚部完全由当前状态和读取符号的虚部通过布尔函数决定}，而不能由外部直接赋值。这保证了虚部标记的``封闭性''：一旦计算进入基域（虚部为$0$），它就无法``凭空''产生新的不可公度元；反之，不可公度元的传播受限于由读取符号虚部和当前状态决定的布尔函数。这一封闭性保证了不可公度性确实由操作语义所内建，而非由外部随意注入——这正是对经典NP定义中``语义外包''的系统性修正。
	
	$M$的计算是一棵配置树。如果至少有一个分支到达$q_{\text{accept}}$，则接受输入；如果所有分支都到达$q_{\text{reject}}$，则拒绝输入。
	
	\subsection{双带视角与代数抽象}
	
	CBTM的纸带可以可视化为两条同步的布尔纸带：
	\begin{itemize}
		\item \textbf{实带}：包含每个单元格的$\operatorname{Re}(\tau)$，承载$\mathbb{F}_2$上的布尔信息，执行传统的确定性布尔计算。
		\item \textbf{虚带}：包含每个单元格的$\operatorname{Im}(\tau)$，标记非确定性分支的发生点。
	\end{itemize}
	两条带共享同一带头位置。当$\operatorname{Im}(\tau) = 0$时，计算保持在基域内；当$\operatorname{Im}(\tau) = 1$时，虚带触发分裂，计算利用扩张域结构产生两条不可公度且互为对偶的延续路径。
	
	双带视角不仅提供了直观的视觉分解，更揭示了CBTM与代数数域之间的深层联系。实带上的有限个连续比特可以组合成一个二进制串，该串可被解释为两个整数的比值。同理，虚带上的连续比特也可组合成另一个有理数。因此，在任意计算时刻，纸带上的一段连续单元格所承载的信息可被抽象为形如
	\[
	a + b \cdot \alpha, \quad a, b \in \mathbb{Q}
	\]
	的代数元，其中系数$a$由实带比特串编码，系数$b$由虚带比特串编码。这一抽象将CBTM的局部存储单元从孤立的布尔值提升为有理数域$\mathbb{Q}$的二次扩域$\mathbb{Q}(\alpha)$中的元素。当$b=0$时，存储退化为普通有理数；当$b \neq 0$时，存储携带了不可公度元$\alpha$的代数信息。这一视角为后续虚部验证机（IVM）中代数维度的度量提供了理论基础。
	
	\textbf{修正经典错误：代码与数据的严格分离，不可公度性的保留。} 在双带视角下，CBTM恢复了图灵机架构中代码与数据分离的基本原则。实部是数据层，承载确定性计算的内容；虚部是控制层，承载非确定性控制信息。见证$c$的每一个比特不再是被拷贝到实部的数据，而是常驻于虚部的函数指针地址。这修正了经典NP定义中将函数指针强制转换为数据的第一重概念错误。
	
	更为重要的是，虚部中的不同生成元（如$\sqrt{2}$和$\sqrt{3}$）在基域上不可公度，这一代数属性\textbf{被操作语义内在地保留}——CBTM通过分支触发公理强制不同分支步骤绑定不同的生成元，从而保证了非确定性选择的独立性。这与经典NP定义形成鲜明对比：后者在将见证拷贝到实部的过程中，不可公度性已被完全抹除。
	
	需要注意的是，$a$和$b$本质上是不同的。$a$是纯粹的数据，而$b$则可以理解为函数指针的地址——它标记了当非确定性分支发生时，后续计算应当调用哪个函数族。经典NP定义正是将这一函数指针强制转换为数据，从而同时造成了身份扭曲和不可公度性丢失。
	
	\subsection{函数指针的类比：虚部作为控制标记}
	
	当虚部$b=1$时，CBTM面对的不是两个简单的``选择''，而是两个\textbf{函数指针}。$\alpha$和$\beta$对应于两个不同的函数族，各自携带不同的代数身份，且互为对偶。机器分裂为两个分支，每个分支各自调用自己的函数指针。确定性操作（虚部为0的步骤）只能访问实部数据，无法感知指针的身份。因此，不存在通用算法能将一条函数族的执行轨迹转化为另一条，这正是不可公度路径独立性的体现。
	
	这一类比揭示了CBTM设计的深层含义。在经典NTM中，``有多个后继''仅是一个语法事实——为什么这些后继构成真正的非确定性选择？标准定义从未回答。但在CBTM中，虚部$b=1$不仅是一个标记，它标记的是\textbf{不可公度元的存在与否}，而分支的两条路径分别绑定于对偶生成元。从而将非确定性选择的独立性和对偶性从语法约定提升为代数事实。
	
	\subsection{NTM的二元标准化}
	
	在建立CBTM与经典模型的等价性之前，我们首先证明一个辅助引理：任何经典的$k$-way分支NTM都可以被转化为一个$2$-way分支的NTM，且转化保持多项式运行时间。这一引理保证了CBTM的二元分支结构在表达能力上没有损失。
	
	\begin{lemma}[NTM的二元标准化]\label{lem:binary-ntm}
		设$N$为一台多项式时间非确定性图灵机，其转移关系在每一步至多产生$k \ge 2$个后继选择。则存在一台多项式时间非确定性图灵机$N'$，满足：
		\begin{enumerate}
			\item $N'$在每一步至多产生$2$个后继选择；
			\item $\mathcal{L}(N') = \mathcal{L}(N)$；
			\item $N'$的运行时间至多为$N$运行时间的常数倍。
		\end{enumerate}
	\end{lemma}
	\begin{proof}
		构造思路是将$N$的一步$k$-way分支展开为一棵深度为$\lceil \log_2 k \rceil$的二叉树。引入至多$\lceil \log_2 k \rceil$个新的中间状态，每个中间状态产生两个后继，分别对应于原选择的前半部分和后半部分。经过至多$\lceil \log_2 k \rceil$步后，每个原选择被隔离到唯一的叶节点。由于$k$是常数，$\lceil \log_2 k \rceil = O(1)$，模拟时间保持多项式界。详细构造见附录A。
	\end{proof}
	
	引理\ref{lem:binary-ntm}表明，在不损失计算能力的前提下，我们可以始终假设经典NTM在每个非确定性步骤恰好有两个选择。这恰好与$\mathbb{F}_4/\mathbb{F}_2$扩张度为$2$的代数事实完美对应，为CBTM与经典模型之间的等价性奠定了组合基础。
	
	\subsection{参数化等价定理}
	
	\begin{theorem}[参数化等价定理]\label{thm:equivalence}
		以下多项式等价性成立：
		\begin{enumerate}
			\item $\text{CBTM}|_0 \equiv_{\text{poly}} \text{DTM}$。因此$\mathbf{P}_{cb} = \mathbf{P}$。
			\item $\text{CBTM} \equiv_{\text{poly}} \text{NTM}$。因此$\mathbf{NP}_{cb} = \mathbf{NP}$。
		\end{enumerate}
	\end{theorem}
	\begin{proof}
		(1) $\text{CBTM}|_0 \to \text{DTM}$：由于所有虚部为$0$，每步转移是确定性的。模拟器逐步骤复制CBTM的转移，线性时间完成。$\text{DTM} \to \text{CBTM}|_0$：将经典DTM的符号编码为$\mathbb{F}_2$中对应元素（虚部设为$0$），转移函数保持原逻辑不变。
		
		(2) $\text{CBTM} \to \text{NTM}$：模拟器将CBTM的每个格局映射为NTM格局，分支时NTM模拟CBTM的两个后继选择。每一步模拟开销为常数时间。$\text{NTM} \to \text{CBTM}$：由引理\ref{lem:binary-ntm}，任何经典NTM可被标准化为二元分支形式。CBTM在遇到$\operatorname{Im}(\tau)=1$时恰好产生两个分支，分别绑定$\alpha$和$\beta$。标准化后的NTM的状态和纸带更新直接编码入CBTM的转移关系中。详细证明见附录B。
	\end{proof}
	
	\begin{remark}[参数化等价定理的双重意义]
		参数化等价定理确立了CBTM与经典计算模型在\textbf{语言识别能力}上的等价性（外延等价）。CBTM并未引入超计算能力。然而，这一等价性仅在外延层面成立——在操作语义层面，CBTM拥有经典模型所不具备的内涵：
		\begin{enumerate}
			\item \textbf{不可公度性内建}：CBTM将不可公度性作为操作语义的内在属性，而非外部语义赋值。经典NTM的``多个后继''仅是语法约定；CBTM的``分叉''是由虚部标记的不可公度元触发的代数事实。经典NP定义在编码过程中丢失了不可公度性；CBTM将其恢复并内建。
			\item \textbf{对偶性显式保留}：CBTM的两条分支分别绑定于$\alpha$和$\beta$，两者互为对偶，由伽罗瓦自同构精确映射。经典NP定义将分支选择编码为比特串并用存在量词消除对偶性；CBTM将这一被消除的结构重新显式化。
			\item \textbf{消除拷贝操作}：CBTM从根本上消除了经典NP定义中隐含的不可实现操作——函数指针地址常驻虚部，分支直接由虚部标记触发，无需将控制信息拷贝到实部。这修正了经典NP定义的操作不自足性。
			\item \textbf{语义内建}：非确定性不再由外部量词``$\exists c$''来定义，而是由机器内部的分支触发公理来执行。
		\end{enumerate}
		经典NTM对CBTM的模拟虽然保持了计算能力，但模拟过程中必须将虚部所携带的代数语义（不可公度性、对偶性）\textbf{编码}为状态信息或额外的纸带符号。CBTM中由代数结构直接触发的分支，在经典模拟中变成了纯粹的语法操作——不同分支之间的代数独立性和对偶性在经典模型的语法层面完全不可见。这种模拟\textbf{丢弃}了CBTM操作语义中的代数结构，而丢弃的代价\textbf{可能是指数级别的复杂度膨胀}。
		
		因此，参数化定理不仅确立了新框架与经典框架的相关性，更揭示了经典框架的语义贫乏和操作不自足。CBTM不是经典NTM的``换壳''，而是为经典NTM\textbf{补充了其长期缺失的语义前提}——不可公度性和对偶性——并\textbf{修正了经典NP定义中的操作不自足和概念扭曲}。这正是``参数化''一词的深层含义：经典NTM的代码与CBTM的代码本质相同，唯一的区别是CBTM将经典模型缺失的语义参数（虚部$b$）显式化并赋予了代数含义，同时消除了经典定义对不可能操作的依赖。
	\end{remark}
	
	\subsection{参数化猜想：$b$是否有意义？}
	
	CBTM通过单一的二进制参数$b$刻画了确定性与非确定性的边界。一个根本性的理论问题是：这个边界是否\textbf{本质}？即，虚部$b$是否为非确定性计算提供了不可消除的代数资源？
	
	\begin{definition}[$b$的有意义性]\label{def:b-meaningful}
		虚部$b$称为\textbf{有意义的}，如果存在语言$L \in \mathbf{NP}$满足：对于任何正确判定$L$的CBTM $M$，存在输入$x$使得在$M$处理$x$的过程中必须使用至少一个满足$\operatorname{Im}(\tau) = 1$的纸带符号。
	\end{definition}
	
	\begin{definition}[$b$的可消除性]\label{def:b-eliminable}
		虚部$b$称为\textbf{可消除的}，如果对于每个$L \in \mathbf{NP}$，都存在一台判定$L$的CBTM，其在所有输入上的所有符号均满足$\operatorname{Im}(\tau)=0$。
	\end{definition}
	
	\begin{conjecture}[参数化猜想]\label{conj:parametrization}
		虚部$b$是有意义的。等价地，存在$\mathbf{NP}$语言使得$b$不可消除。
	\end{conjecture}
	
	\begin{theorem}[参数化等价]\label{thm:param-equivalence}
		参数化猜想$\iff \mathbf{P} \neq \mathbf{NP}$。
	\end{theorem}
	\begin{proof}
		($\Rightarrow$) 设参数化猜想成立，则存在$L \in \mathbf{NP}$不能被任何$\text{CBTM}|_0$判定。由定理\ref{thm:equivalence}(1)，$L \notin \mathbf{P}$，故$\mathbf{P} \neq \mathbf{NP}$。
		
		($\Leftarrow$) 设$\mathbf{P} \neq \mathbf{NP}$。若$b$可消除，则对任意$L \in \mathbf{NP}$存在$\text{CBTM}|_0$判定$L$，从而$L \in \mathbf{P}$，与$\mathbf{P} \neq \mathbf{NP}$矛盾。故$b$不可消除。
	\end{proof}
	
	定理\ref{thm:param-equivalence}将Cook和Levin提出的经典问题\textbf{等价转化}为一个关于虚部$b$的代数问题：域扩张$\mathbb{F}_4/\mathbb{F}_2$是否为非确定性计算提供了本质的代数资源？参数化猜想的最终解决——即$b$是否有意义——依赖于后续论文中虚部验证机（IVM）的动态维度追踪机制和本质维度$\kappa(L)$的严格论证。我们将在第5.8节讨论CBTM静态框架的局限性，并指出走向动态维度追踪的必要性。
	
	
	\section{实布尔图灵机与生成元无关性}
	
	CBTM利用抽象生成元$\alpha$刻画非确定性计算，但$\alpha$的抽象性使其缺乏直观的几何意义。本节将$\alpha$具体化为代数数$\sqrt{2}$，引入实布尔图灵机（RBTM）和双带视角，并证明生成元无关性定理：计算能力不依赖于生成元的具体选择。在此基础上，我们揭示伽罗瓦自同构的操作本质——路径切换——并据此提出敏感性猜想，将P vs.\ NP问题转化为一个定性结构判据。
	
	\subsection{从$\alpha$到$\sqrt{2}$：RBTM的形式化定义}
	
	CBTM中的生成元$\alpha$满足$\alpha^2=\alpha+1$，是一个抽象符号。为了获得直观的几何意义，我们将其替换为具体的代数数$\sqrt{2}$。$\sqrt{2}$与有理数域$\mathbb{Q}$不可公度，这正是作为``新维度''所需的性质。而且，$\sqrt{2}$与其对偶$-\sqrt{2}$满足完全相同的代数方程$x^2-2=0$，由伽罗瓦自同构精确映射。于是，每个符号可以表示为
	\[
	\tau = a + b\sqrt{2},\quad a,b\in\{0,1\}.
	\]
	实部$a$编码传统布尔值，虚系数$b$指示是否涉及$\sqrt{2}$。这一替换不仅是符号层面的变化，更重要的是将CBTM的离散有限域结构自然地推广到无限域，为后续的双带视角和生成元无关性提供了直观的几何基础。
	
	\begin{definition}[RBTM]\label{def:rbtm}
		一台实布尔图灵机定义为一个七元组
		\[
		M_{\mathrm{RBTM}} = (Q,\Sigma,\Gamma,\delta,q_0,F,\epsilon),
		\]
		其中：
		\begin{itemize}
			\item $Q$：有限状态集；
			\item $\Sigma \subseteq \mathbb{Q}(\sqrt{2})$：输入字母表，每个符号形如$a+b\sqrt{2}$，$a,b\in\{0,1\}$；
			\item $\Gamma = \Sigma \cup \{\#\}$，$\#$为空白符号，其投影定义为$\mathfrak{Re}(\#)=\mathfrak{Im}(\#)=0$；
			\item $\delta: Q\times\Gamma \to \mathcal{P}(Q\times\Gamma\times\{L,R\})$：转移函数，满足分支触发公理和投影约束公理（将$\alpha$替换为$\sqrt{2}$）；
			\item $q_0\in Q$：初始状态；
			\item $F\subseteq Q$：接受状态集；
			\item $\epsilon\in\mathbb{Q}^+$：分支阈值（通常取$\epsilon=0.5$）。
		\end{itemize}
	\end{definition}
	
	RBTM与CBTM的计算等价性由以下定理保证。该定理证明，将抽象生成元替换为具体代数数不改变计算能力。
	
	\begin{theorem}[RBTM与CBTM的等价性]\label{thm:rbtm-cbtm-equiv}
		存在多项式时间的相互模拟，使得$\mathcal{L}(\mathrm{RBTM})=\mathcal{L}(\mathrm{CBTM})$，且复杂度类相等：$\mathbf{P}_{rb}=\mathbf{P}_{cb}=\mathbf{P}$，$\mathbf{NP}_{rb}=\mathbf{NP}_{cb}=\mathbf{NP}$。
	\end{theorem}
	\begin{proof}
		定义符号映射$\psi:\mathcal{CB}\to\mathbb{Q}(\sqrt{2})$为：
		\[
		\psi(0)=0,\quad \psi(1)=1,\quad \psi(\alpha)=\sqrt{2},\quad \psi(\beta)=1+\sqrt{2}.
		\]
		易验证$\psi$保持投影值：$\mathfrak{Re}_{\mathrm{CBTM}}(\tau)=\mathfrak{Re}_{\mathrm{RBTM}}(\psi(\tau))$，$\mathfrak{Im}_{\mathrm{CBTM}}(\tau)=\mathfrak{Im}_{\mathrm{RBTM}}(\psi(\tau))$。将CBTM的每个格局通过$\psi$逐符号映射，即得到RBTM的格局；反之用逆映射将RBTM格局映回CBTM。由于转移函数由投影值决定，该映射保持转移关系，每一步开销为$O(1)$。故二者多项式时间等价。
	\end{proof}
	
	\subsection{双带视角：让维度可见}
	
	RBTM的纸带可以物理分解为两条平行的布尔纸带：
	\begin{itemize}
		\item \textbf{实带}：在每个位置存储实系数$a$，执行传统的确定性布尔计算；
		\item \textbf{虚带}：在每个位置存储虚系数$b$，标记``新维度''出现的位置。
	\end{itemize}
	两条带共享同一带头位置。这种分解建立了自然同构$\mathcal{CB}\cong\{0,1\}\times\{0,1\}$，并且保持结构。
	
	双带视角提供了强有力的直觉：实带执行传统的确定性布尔计算；虚带控制非确定性——当虚带上某个位置为1时，意味着``此处嵌入了一个新维度''，机器必须分叉，同时探索包含和排除该维度的两种可能性。虚带上的每一个``1''都直观地标志着新维度的存在，使``维度''概念首次变得可视化。这一可视化不仅具有直观价值，更为后续虚部验证机（IVM）中动态维度追踪奠定了物理基础。
	
	考虑一台RBTM处理单个符号$\sqrt{2}$（即$\tau=0+1\cdot\sqrt{2}$）。在双带视角下，该位置实带为0，虚带为1。虚部为1触发分支：左分支对应``包含$\sqrt{2}$''，写入$1+\sqrt{2}$（实部1，虚部1）；右分支对应``排除$\sqrt{2}$''，写入$0$（实部0，虚部0）。值得注意的是，左分支绑定了生成元$+\sqrt{2}$，右分支的排除并非退回到基域（该地址仍然被分配），而是绑定其对偶$-\sqrt{2}$。这一示例展示了虚部如何通过布尔标记直接控制计算的分叉行为，以及对偶性如何被保留在操作语义中。
	
	\subsection{域非同构问题的消解}
	
	CBTM的代数基础是有限域$\operatorname{GF}(4)$（特征2），而RBTM基于无限域$\mathbb{Q}(\sqrt{2})$（特征0）。作为代数结构，这两个域一个是有限的，一个是无限的，并不同构。然而，在双带视角下，每个符号被分解为两个布尔分量（实部和虚部），计算行为完全由这两个布尔分量的演化决定，而与生成元的具体代数性质无关。
	
	因此，尽管$\operatorname{GF}(4)$与$\mathbb{Q}(\sqrt{2})$在结构上存在差异，它们在布尔投影层所诱导的自动机却是\textbf{自动机同构}的。这是生成元无关性定理的一个推论：不同的代数表示可以产生完全相同的计算行为。双带视角将代数符号``投影''到布尔层，消除了底层域结构的差异，使我们能够专注于计算的本质——新维度的有无，而非维度的代数身份。
	
	更重要的是，这一视角揭示出\textbf{非确定性的代数本质——不可公度性与伽罗瓦群作用（即对偶性）——在比特层面被不变地保留下来，超越了有限域与无限域之间的鸿沟}。无论底层是特征2的有限域$\operatorname{GF}(4)$，还是特征0的无限域$\mathbb{Q}(\sqrt{2})$，分支选择序列的翻转操作（对应于伽罗瓦群在路径上的作用，即对偶路径切换）在布尔投影层完全一致。计算语义的同一性因此超越了代数结构的差异，这正是生成元无关性定理的深层内涵。
	
	\subsection{生成元无关性定理}
	
	生成元无关性定理是本节的核心技术贡献。令$\gamma$和$\delta$为任意两个非有理代数数（例如$\sqrt{2}$与$\sqrt{3}$，或$\sqrt{2}$与$i$）。分别以$\gamma$和$\delta$为生成元构造RBTM $M_\gamma$和$M_\delta$。我们要证明：$M_\gamma$与$M_\delta$在计算上完全等价。
	
	核心观察在于：分支触发仅取决于虚系数的布尔值（$\mathfrak{Im}(\tau)=b$），与生成元的具体数值无关。无论生成元是$\sqrt{2}$还是$i$，只要虚系数$b=1$，机器就必须分叉，且分叉的两条路径分别绑定于对偶生成元。因此，两台机器的任何行为差异只能源于符号上的代数运算，但投影算子仅提取有理系数$a,b$，这些系数在生成元替换下保持不变。
	
	\begin{lemma}[投影值的一致性]\label{lem:proj-consistency}
		对任意符号$\tau=a+b\gamma$（$a,b\in\mathbb{Q}$），其在生成元$\gamma$下的投影为$\mathfrak{Re}_\gamma(\tau)=a$，$\mathfrak{Im}_\gamma(\tau)=b$。若将$\gamma$替换为另一生成元$\delta$，得到$\tau'=a+b\delta$，则$\mathfrak{Re}_\delta(\tau')=a$，$\mathfrak{Im}_\delta(\tau')=b$。即投影值仅依赖于有理系数$a,b$，与生成元无关。
	\end{lemma}
	\begin{proof}
		由投影算子的定义直接可得。
	\end{proof}
	
	\begin{definition}[自动机同构]\label{def:automaton-isomorphism}
		两台自动机$M_1=(Q_1,\Sigma_1,\delta_1,q_{01},F_1)$与$M_2=(Q_2,\Sigma_2,\delta_2,q_{02},F_2)$称为同构（记作$M_1\cong M_2$），如果存在双射$\phi:Q_1\to Q_2$满足：
		\begin{enumerate}
			\item $\phi(q_{01})=q_{02}$；
			\item $\phi(F_1)=F_2$；
			\item 对任意$q\in Q_1$和$\tau\in\Sigma_1$，$\phi(\delta_1(q,\tau))=\delta_2(\phi(q),\tau)$（这里我们将$\delta$的值域解释为状态集，忽略写入符号和带头移动，因为这些可以编码进状态；严格处理需要格局同构，但为简化，我们专注于状态转移的对应；纸带符号的同构由投影值的一致性保证）。
		\end{enumerate}
	\end{definition}
	
	\begin{theorem}[生成元无关性]\label{thm:generator-independence}
		设$\gamma$和$\delta$为任意两个非有理代数数。则基于$\gamma$的RBTM $M_\gamma$与基于$\delta$的RBTM $M_\delta$自动机同构。特别地，它们识别相同的语言类，且具有相同的复杂度类。
	\end{theorem}
	\begin{proof}
		首先，构造字母表间的双射。记$\Gamma_\gamma = \{0,1,\gamma,1+\gamma\}$，$\Gamma_\delta = \{0,1,\delta,1+\delta\}$。定义映射$\psi:\Gamma_\gamma\to\Gamma_\delta$为
		\[
		\psi(a+b\gamma)=a+b\delta,\qquad a,b\in\{0,1\}.
		\]
		即$\psi(0)=0,\;\psi(1)=1,\;\psi(\gamma)=\delta,\;\psi(1+\gamma)=1+\delta$。显然$\psi$是双射且保持投影值：
		\[
		\mathfrak{Re}_\delta(\psi(\tau)) = a = \mathfrak{Re}_\gamma(\tau),\qquad
		\mathfrak{Im}_\delta(\psi(\tau)) = b = \mathfrak{Im}_\gamma(\tau).
		\]
		
		状态映射取恒等映射$\iota:Q_\gamma\to Q_\delta$（两台机器的状态集可视为相同，因为它们仅依赖于控制逻辑，与生成元的具体值无关）。现验证在$(\iota,\psi)$下转移关系得以保持。
		
		考虑$M_\gamma$中的任意转移$(q',\gamma_w,d)\in\delta_\gamma(q,\tau)$，其中$\tau=a+b\gamma$。根据投影约束公理，存在布尔函数$f_{\mathfrak{Re}},f_{\mathfrak{Im}}$使得写入符号$\gamma_w$满足：
		\[
		\gamma_w = f_{\mathfrak{Re}}(q,\mathfrak{Re}_\gamma(\tau)) + \gamma\cdot f_{\mathfrak{Im}}(q,\mathfrak{Im}_\gamma(\tau)).
		\]
		在$M_\delta$中，读取符号为$\psi(\tau)=a+b\delta$，其投影值与$\tau$相同。由同样的布尔函数$f_{\mathfrak{Re}},f_{\mathfrak{Im}}$，根据投影约束公理，$M_\delta$在状态$\iota(q)=q$下必有一个转移$(q',\delta_w,d)$使得
		\[
		\delta_w = f_{\mathfrak{Re}}(q,\mathfrak{Re}_\delta(\psi(\tau))) + \delta\cdot f_{\mathfrak{Im}}(q,\mathfrak{Im}_\delta(\psi(\tau))) = f_{\mathfrak{Re}}(q,a) + \delta\cdot f_{\mathfrak{Im}}(q,b).
		\]
		注意到$\psi(\gamma_w)=f_{\mathfrak{Re}}(q,a)+\delta\cdot f_{\mathfrak{Im}}(q,b)=\delta_w$，故$(q',\psi(\gamma_w),d)\in\delta_\delta(q,\psi(\tau))$。反之，$M_\delta$中的每个转移也对应$M_\gamma$中唯一的转移。因此转移关系在$(\iota,\psi)$下得以保持。
		
		初始状态和接受状态显然满足$\iota(q_{0\gamma})=q_{0\delta}$，$\iota(F_\gamma)=F_\delta$。故$M_\gamma\cong M_\delta$。自动机同构保持语言识别和复杂度类，结论得证。详细补充见附录C。
	\end{proof}
	
	生成元无关性定理具有深远的理论意义。第一，它确认了非确定性计算的代数本质不依赖于生成元的个体身份，而仅依赖于``存在与基域不可公度的元素''这一事实，以及生成元之间的对偶关系。第二，它赋予了我们在后续工作中自由选择最方便生成元体系的权力（例如，在IVM中使用素数平方根序列$\sqrt{p_1},\sqrt{p_2},\dots$标记不同的分支步骤），而不必担心计算能力发生改变。第三，结合双带视角，该定理表明尽管有限域$\operatorname{GF}(4)$和无限域$\mathbb{Q}(\sqrt{2})$在代数上不同构，它们在布尔投影层所诱导的自动机却是同构的。计算语义的同一性——包括不可公度性和对偶性——可以超越底层域结构的差异。
	
	\subsection{伽罗瓦自同构的操作本质：路径切换}
	
	生成元无关性定理揭示了一个更深层的洞察：在比特投影层，伽罗瓦自同构$\sigma$的作用完全一致，与生成元无关。那么，$\sigma$在计算树上究竟做了什么？
	
	\textbf{答案异常简洁：$\sigma$将每条路径的所有非确定性选择全部翻转。}
	
	形式化地，设一条计算路径$\pi$包含$m$个非确定性选择点，其分支选择序列为$\mathbf{d} = d_1 d_2 \ldots d_m \in \{0,1\}^m$。在CBTM中，$d_i=0$对应于选择$\alpha$路径，$d_i=1$对应于选择$\beta$路径（或反之，取决于具体的编码约定）。伽罗瓦自同构$\sigma: \alpha \leftrightarrow \beta$在计算树上的作用，正是将每条路径$\pi$全局地切换为其对偶路径$\sigma(\pi)$：
	\[
	\sigma(\pi) = \text{将 } \mathbf{d} \text{ 替换为 } \mathbf{d} \oplus \mathbf{1}_m \text{ 所得到的路径},
	\]
	其中$\oplus$表示逐位异或，$\mathbf{1}_m$是长度为$m$的全1向量。这一操作将路径上的每一个非确定性选择全部翻转：原本选择$\alpha$的变为选择$\beta$，原本选择$\beta$的变为选择$\alpha$。
	
	这一揭示的重要意义在于：伽罗瓦自同构不再是抽象的代数对象，而是获得了明确的操作性解释——它就是对计算树中所有非确定性选择的一次``全局翻转''。这是对偶性在操作层面的直接体现：非确定性选择的两个选项不是任意的，而是通过伽罗瓦群作用精确对偶的。在经典NP定义中，这一对偶性被存在量词消除；在CBTM中，它被显式保留并获得了操作性的意义。在虚部验证机（IVM）的动态扩张下，每次非确定性分支被分配一个独立的生成元，伽罗瓦群从单一的$\mathbb{Z}/2\mathbb{Z}$扩展为$(\mathbb{Z}/2\mathbb{Z})^m$。这使得我们可以独立地翻转任意非确定性选择点处的选择，从而将全局路径切换分解为对单个选择点的独立操作。
	
	\subsection{单点敏感性与路径切换判据}
	
	基于上述操作化，我们可以定义更为精细的敏感性概念。
	
	\begin{definition}[单点敏感性]\label{def:pointwise-sensitivity}
		称机器$M$在输入$x$上\textbf{对第$i$个非确定性选择点敏感}，如果存在一条接受路径$\pi$，使得单独翻转$\pi$上第$i$个选择后（即将$d_i$替换为$1-d_i$，而保持其他选择不变），所得路径$\pi^{\oplus i}$不再是接受路径。若对某个输入$x$，存在某个非确定性选择点$i$使得$M$敏感，则称$M$\textbf{单点敏感}。若$M$在所有输入上对所有选择点均不敏感，则称$M$\textbf{完全不敏感}。
	\end{definition}
	
	这一翻转操作在代数上对应于将生成元的选择从$+\sqrt{p_i}$切换为其对偶$-\sqrt{p_i}$（或反之），即伽罗瓦自同构在单个非确定性选择点上的作用。因此，单点敏感性本质上是对偶性在操作层面的直接体现。
	
	单点敏感性与全局敏感性（伽罗瓦自同构同时翻转所有选择）的关系是：若机器是全局敏感的，则它至少在某一个选择点上是单点敏感的；反之，若机器是完全不敏感的（对所有单点翻转均不改变结果），则它也是全局不敏感的。因此，单点敏感性是比全局敏感性更精细的刻画，它直接对应每个非确定性选择是否本质。在后续的IVM框架中，本质维度$\kappa(L)$正是通过度量被激活的生成元（对应本质的选择点）来量化这一概念。
	
	\begin{theorem}[P vs NP的路径切换判据]\label{thm:path-switch}
		\begin{enumerate}
			\item 若语言$L$可由一台对所有非确定性选择点\textbf{完全不敏感}的CBTM/RBTM判定，则$L \in \mathsf{P}$。
			\item 若$\mathsf{P} \neq \mathsf{NP}$，则对任何$L \in \mathsf{NP} \setminus \mathsf{P}$，任何接受$L$的CBTM/RBTM在某些输入上必然存在某个非确定性选择点，使得机器在该点上\textbf{单点敏感}。特别地，NP完全语言在任何接受它的机器上都具有\textbf{不可消除的单点敏感性}。
		\end{enumerate}
	\end{theorem}
	\begin{proof}
		(1) 设$M$是一台对所有非确定性选择点完全不敏感的CBTM/RBTM，判定语言$L$。对于任意输入$x$，考虑$M$的计算树。设树中非确定性选择点总数为$m$。定义函数$f: \{0,1\}^m \to \{0,1\}$，$f(\mathbf{d}) = 1$当且仅当分支选择序列为$\mathbf{d}$的路径是接受路径。
		
		完全不敏感性意味着：对于任意$\mathbf{d} \in \{0,1\}^m$和任意$i \in \{1,\dots,m\}$，有$f(\mathbf{d}) = f(\mathbf{d} \oplus \mathbf{e}_i)$，其中$\mathbf{e}_i$是仅第$i$位为1的单位向量。通过逐位翻转，可以从任意$\mathbf{d}$到达任意$\mathbf{d}'$，因此$f$是常数函数。
		
		既然$f$是常数，机器只需取任意一条固定路径（如全0序列对应的路径）即可得到正确的接受/拒绝结果。这条固定路径可以在确定性多项式时间内模拟：在每一步非确定性选择点处，固定选择分支0。由于$f$是常数，这条路径的接受状态与任何其他路径相同。因此$L \in \mathsf{P}$。
		
		(2) 若$\mathsf{P} \neq \mathsf{NP}$，设$L \in \mathsf{NP} \setminus \mathsf{P}$。假设存在一台接受$L$的机器$M$在所有输入上对所有选择点均完全不敏感，则由(1)知$L \in \mathsf{P}$，矛盾。故任何接受$L$的机器必然在某些输入上存在至少一个单点敏感的选择点。对于NP完全语言，若某台接受它的机器可被修改为完全不敏感，则可通过(1)的构造得到确定性多项式时间算法，从而$\mathsf{P} = \mathsf{NP}$。因此NP完全语言的单点敏感性是\textbf{不可消除的}。
	\end{proof}
	
	定理\ref{thm:path-switch}是本文的核心贡献之一。它将P vs NP问题从``计算需要多少时间''的定量资源分析，转化为一个简洁的定性结构判据：
	\begin{enumerate}
		\item 若在所有非确定性选择点上单独翻转从不改变计算结果，则非确定性是冗余的，$\mathsf{P} = \mathsf{NP}$；
		\item 若存在某个非确定性选择点，单独翻转必然改变某些计算结果，则非确定性是本质的，$\mathsf{P} \neq \mathsf{NP}$。
	\end{enumerate}
	
	这一刻画的意义在于：它将P与NP的区分从定量问题转化为定性问题，这一区分建立在对偶性的操作化之上——单点敏感性正是对偶性在单个选择点上的体现。伽罗瓦敏感性在比特投影层是不变的——无论底层是有限域$\operatorname{GF}(4)$还是无限域$\mathbb{Q}(\sqrt{2})$，翻转单个分支选择的操作完全相同。因此，复杂性的分界线\textbf{超越}了具体的代数实现，成为纯粹的结构性事实。
	
	\subsection{敏感性猜想及其与$\mathbf{P}\neq\mathbf{NP}$的等价性}
	
	基于定理\ref{thm:path-switch}，我们将P与NP的区分问题提炼为如下猜想。
	
	\begin{conjecture}[敏感性猜想]\label{conj:sensitivity}
		存在一个语言$L\in\mathsf{NP}$，使得任何判定$L$的多项式时间CBTM（或RBTM）都必然在某个输入上对至少一个非确定性选择点单点敏感。换言之，不存在一台完全不敏感的机器能够判定某个NP完全语言。
	\end{conjecture}
	
	\begin{theorem}[等价性]\label{thm:sensitivity-equivalence}
		敏感性猜想等价于$\mathsf{P}\neq\mathsf{NP}$。
	\end{theorem}
	\begin{proof}
		由定理\ref{thm:path-switch}直接可得：(1) 若$\mathsf{P}\neq\mathsf{NP}$，则任何NP完全语言均满足敏感性猜想的条件；(2) 反之，若存在这样的语言，则其不能被任何完全不敏感的机器判定，从而由定理\ref{thm:path-switch}(1)知其不属于$\mathsf{P}$，故$\mathsf{P}\neq\mathsf{NP}$。
	\end{proof}
	
	\textbf{敏感性猜想}将P vs NP问题转化为一个关于计算路径对单点翻转的响应性质的命题：非确定性计算的力量是否必然导致存在某个本质的非确定性选择点，单独翻转其选择就会改变接受状态？这一本质选择点正是生成元对偶性在操作层面的直接体现——如果每个选择点的对偶翻转都不影响结果，那么对偶性就是冗余的，非确定性可以被确定性模拟消除；反之，如果存在某个选择点其对偶翻转必然改变结果，那么对偶性就是本质的，非确定性无法被消除。猜想断言这种本质选择点的存在是NP完全语言不可消除的内在属性。在后续的IVM框架中，每一个本质的非确定性选择点恰好对应一个被激活的生成元，因此本质维度$\kappa(L)$恰好度量了本质选择点的数量。敏感性猜想的成立将直接导出$\kappa(L) = \Omega(n)$的下界，从而完成$\mathsf{P} \neq \mathsf{NP}$的证明。
	
	\subsection{生成元提取算子与静态局限性}
	
	尽管RBTM的双带视角使``新维度''可视，但所有分支共享同一个生成元$\sqrt{2}$，无法区分不同分支步骤所引入的维度。为了精确量化计算过程中的代数资源消耗，我们需要一个能从每个符号中提取生成元身份的算子。
	
	\begin{definition}[生成元提取算子]\label{def:generator-extraction}
		对于符号$\tau = a+b\sqrt{2}\in\mathbb{Q}(\sqrt{2})$，定义生成元提取算子$\mathfrak{G}:\Gamma\to\{\sqrt{2},1,0\}$为：
		\[
		\mathfrak{G}(\tau) = \begin{cases}
			\sqrt{2} & \text{若 } b\neq 0,\\
			1 & \text{若 } b=0\text{ 且 } a\neq 0,\\
			0 & \text{若 } a=b=0.
		\end{cases}
		\]
	\end{definition}
	
	然而，在静态RBTM框架内，$\mathfrak{G}(\tau)$只能返回$\sqrt{2}$、$1$或$0$，无法区分多个不同分支步骤所贡献的维度，因为所有分支步骤指向同一个$\sqrt{2}$。这就是静态框架的\textbf{根本局限性}：它只能告诉我们``是否涉及新维度''，却无法回答``涉及了多少个不同的新维度''。对于后续需要量化复杂度的工作（例如定义本质维度$\kappa(L)$），我们需要一个能够动态引入\textbf{多个线性无关生成元}的框架。
	
	具体而言，CBTM/RBTM虽然完美地刻画了单次分支的独立性（不可公度性）和对偶性（伽罗瓦自同构），但无法区分``发生一次分支''与``发生了$n$次分支''——多次分支的代数维度会坍缩回常数$2$，而多次分支所引入的不同生成元的对偶关系也被压缩为同一对对偶元。这一局限正是引入动态的虚部验证机（IVM）的根本动机：IVM通过为每次分支分配独立的素数平方根生成元，将不可公度性升级为线性无关性，同时为每次分支分配独立的对偶对（$\sqrt{p_i}$与$-\sqrt{p_i}$），从而使得代数维度能够随分支次数线性增长，对偶性也得以按分支步骤独立追踪。生成元无关性定理为这一升级提供了理论保证：我们可以自由选择最方便的生成元体系（素数平方根序列），而不改变计算能力。
	
	\section{方法论原则：数计一体与范式转换}
	
	前述各节从元理论诊断（第3节）到新模型的构建（第4-5节），始终贯穿着一条根本性的方法论线索。本节将这条线索明确提炼为``数计一体''原则，并阐明CBTM/RBTM框架如何体现这一原则，以及它所标志的范式转换——特别地，这一框架如何系统性地修正经典NP定义的操作不自足和三重概念错误。
	
	\subsection{数计一体原则的陈述}
	
	\textbf{数计一体}（the unity of mathematics and computation）主张：计算理论应当仿照数学，要求任意算法都应当具有明确的语义，并且语义与语法保持一致。
	
	在数学中，一个表达式不仅有其书写形式（语法），更有其指称的数学对象和满足的代数关系（语义）。数学推理的每一步都同时遵守语法规则和语义约束，两者始终协调一致。数学计算的正确性依赖于这种严格对应。例如，当我们写下$a+b\sqrt{2}$时，这一表达式不仅是一个符号串，还同时携带了它在扩域$\mathbb{Q}(\sqrt{2})$中的代数身份——它与基域元素的不可公度性、它作为二维向量空间中元素的线性结构等。
	
	经典计算理论背离了这一原则。图灵机模型刻意剥离了符号的语义，将计算还原为纯粹的语法操作。这种设计赋予了图灵机通用性，但也使计算模型丧失了``理解''其所处理对象的能力。经典NP定义不得不将非确定性的本质外包给外部量词和外部谓词，并在外包过程中犯下了操作不自足和概念扭曲——函数指针被扭曲为数据（同时丢失不可公度性）、对偶性被量词消除、整个定义依赖于不可实现的拷贝操作——正是这一背离的直接后果。
	
	数计一体原则要求重新统一语法与语义。具体而言，计算模型的设计应当满足以下要求：
	\begin{enumerate}
		\item \textbf{符号自带语义}：纸带上的每一个符号本身就是一个数学对象，而非等待外部赋予含义的裸标记。
		\item \textbf{操作保持语义}：每一步计算操作不仅移动符号，同时也在执行相应的数学运算，且运算结果忠实地反映在符号的语义变化中。
		\item \textbf{语义决定行为}：计算的走向（如是否分支）由符号的语义内容（如是否涉及不可公度元）决定，而非由外部量词或外部谓词事后赋予。
		\item \textbf{消除不可能操作}：计算模型的操作语义必须是自足的——模型不应要求其自身无法执行的操作（如虚部到实部的拷贝）。
		\item \textbf{对偶性显式保留}：非确定性选择的对偶结构应当作为操作语义的内在属性被保留，而非被编码和量词消除。
	\end{enumerate}
	
	\subsection{CBTM/RBTM的参数化设计：对操作不自足与三重错误的系统性修正}
	
	CBTM是数计一体原则的第一个实例，它的设计体现了参数化方法的精髓，并且可以精确地理解为对经典NP定义的操作不自足和三重概念错误的系统性修正。
	
	\textbf{参数化设计：一个不改变代码的一一映射。} CBTM的设计基于一个极为简洁的观察：经典非确定性图灵机（NTM）和确定性图灵机（DTM）的代码结构本质上是相同的，唯一的区别在于NTM在某些格局下允许转移关系中有多个五元组。CBTM对这一结构做了一个巧妙的一一映射，\textbf{完全没有改变TM或NTM的任何代码}：
	\begin{itemize}
		\item 对于DTM的每一个确定性步骤，CBTM将其映射为虚部$b=0$的符号操作——纸带上的内容被解释为基域中的元素，计算在基域内进行，不引入代数扩张。
		\item 对于NTM的每一个非确定性分叉步骤，CBTM将其映射为虚部$b=1$的符号操作——在分叉发生之前，当前纸带格子上的虚部被\textbf{预先标记}为$1$，表示该位置涉及域扩张中的不可公度元。两条分支路径分别绑定于对偶生成元。\textbf{分支直接由虚部标记触发，无需将控制信息拷贝到实部。}
	\end{itemize}
	
	这一映射是双射：给定任何一台DTM或NTM，存在唯一对应的CBTM（代码相同，仅在分叉处加标记）；反之，给定任何一台CBTM，抹去所有虚部标记即可恢复原DTM或NTM（代码相同，仅去掉标记）。\textbf{代码本质相同，虚部仅作为分支标记。}这正是参数化等价定理（定理\ref{thm:equivalence}）在操作层面的对应物——$\mathbf{P}_{cb} = \mathbf{P}$，$\mathbf{NP}_{cb} = \mathbf{NP}$。
	
	\textbf{修正操作不自足：从根本上消除拷贝操作。} 在经典NP定义中，见证$c$从虚部被拷贝到实部，这一操作在经典操作语义中不可实现。在CBTM中，这一问题被从根本上消除：函数指针的地址不再需要拷贝到实部，而是常驻于虚部。非确定性分支由机器内部的分支触发公理直接执行，无需外部观察者介入，无需从控制信息到数据的跨层次转换。
	
	\textbf{修正第一重错误：函数指针常驻虚部，不可公度性得以保留。} 在经典NP定义中，见证$c$被拷贝到实部成为数据，这一转换同时导致不可公度性的丢失。在CBTM中，虚部符号的语义身份是明确的——它是控制信息，标记``此处涉及不可公度元''，而非数据。实部与虚部的严格分离，恢复了代码与数据分离的原则。更为重要的是，虚部中的不同生成元在基域上不可公度，这一代数属性被操作语义内在地保留。
	
	\textbf{修正第二重错误：对偶性显式保留。} 在经典NP定义中，对偶性被编码为比特串并用存在量词消除。在CBTM中，对偶性被显式保留在虚部空间中：分支的两条路径分别绑定于$\alpha$和$\beta$，由伽罗瓦自同构精确映射。对偶性不再是被抹去的隐性结构，而是操作语义的显式代数事实。
	
	\textbf{修正第三重错误：语义内建而非外包。} 在经典NP定义中，非确定性由外部量词``$\exists c$''来断言。在CBTM中，非确定性由机器内部的分支触发公理来执行：当虚部$b=1$时，机器强制分叉——这是由不可公度性触发的代数行为。
	
	正是这一参数化设计，使得非确定性选择的独立性和对偶性获得了正面的代数刻画。两个分支之所以是``真正不同的选择''，不是因为它们对应不同的五元组（语法层面），而是因为它们分别绑定了对偶的不可公度生成元（代数层面）。这种独立性由域扩张的代数结构严格保证，对偶性由伽罗瓦群的作用精确刻画，无法通过基域内的任何运算相互转化。
	
	RBTM通过双带视角和生成元无关性定理进一步推广了这一原则。双带视角将实部（数据）与虚部（控制）分离，使``维度''概念可视化。生成元无关性定理则确保了我们可以自由选择最方便的生成元体系，而不改变计算能力——这为后续IVM中动态分配素数平方根生成元提供了理论保证。
	
	\subsection{与丘奇-图灵论题的关系}
	
	数计一体可以被视为丘奇-图灵论题的一个自然延伸。丘奇-图灵论题断言任何算法可计算的函数都是图灵机可计算的。数计一体则进一步要求：任何数学对象所携带的语义结构（如不可公度性、线性无关性、代数对偶性等）都应当在计算模型的操作语义中获得正面的、内在的表达，且计算模型的操作语义必须是自足的——不应依赖其自身无法执行的操作。
	
	这不是对丘奇-图灵论题的否定，而是对其``语法中立性''预设的修正。丘奇-图灵论题关注的是``能计算什么''的外延等价；数计一体关注的是``如何计算''的内涵充分性和操作自足性。一个计算模型不应仅满足于可计算性的外延等价，而应当追求语义表达的内蕴充分性和操作自足性。CBTM/RBTM框架在保持外延等价的前提下，通过内建不可公度性和保留对偶性丰富了操作语义，并从根本上消除了拷贝操作的必要性，正是这一方法论原则的严格实现。
	
	\subsection{从语法范式到语义范式的转换}
	
	维度退化理论标志着计算理论的一次范式转换。旧范式（经典图灵机）的语言只描述符号的语法操作，无法触及计算背后的代数语义——不可公度性被外包，对偶性被消除，操作不自足被掩盖。新范式（CBTM/IVM）将非确定性的本质——不可公度性及其对偶结构——内建于操作语义，从根本上消除了拷贝操作的必要性，使P与NP的差异从模糊的哲学问题转化为精确的代数对比。
	
	这不是在旧框架内的修补，而是从``语义外包、操作不自足''到``语义内建、操作自足''的根本转变。旧范式下，非确定性被定义为``存在多个后继''（语法约定），选项之间的独立性靠直觉维持，不可公度性在编码中丢失，对偶性被量词消除，定义本身依赖于不可实现的拷贝操作。新范式下，非确定性被定义为``引入不可公度生成元并强制分叉为对偶路径''（代数事实），独立性由生成元的线性无关性严格保证，对偶性由伽罗瓦群的作用显式保留，拷贝操作被彻底消除。
	
	参数化定理正是这一范式转换的数学表达：它证明了经典NTM的代码与CBTM的代码本质相同，区别仅在于后者将前者缺失的语义前提（不可公度性、对偶性）补全了，消除了前者隐含依赖的不可能操作，并修正了经典NP定义在外包过程中造成的概念扭曲。
	
	P vs.\ NP问题由此被转化为一个精确的代数问题：代数数域上的计算是否可以由有理数域在多项式时间内模拟？如果答案是肯定的，那么非确定性计算的代数扩张就是冗余的，可以在基域内被确定性模拟消除，即$\mathbf{P} = \mathbf{NP}$。如果答案是否定的，那么代数扩张提供了本质的计算资源，无法被基域内的仿射变换所模拟，即$\mathbf{P} \neq \mathbf{NP}$。本文已为这一问题提供了模型基础和方法论框架——从不可公度性元定理到CBTM/RBTM的参数化设计，再到敏感性猜想。后续论文将基于此框架对这一问题展开严格分析。
	
	
	\section{结论与展望}
	
	\subsection{本文的主要贡献}
	
	本文从经典图灵机模型的表达局限出发，构建了非确定性计算的代数语义框架，由三个逻辑上紧密关联的部分组成。
	
	\textbf{第一，元理论基石。} 我们基于语法不变性原理，严格证明了不可公度性不能是经典图灵机操作语义的内在属性（定理\ref{thm:main}）。该证明的核心逻辑是：图灵机的转移函数仅依赖符号标识，机器行为在符号置换下同构，因此任何操作语义的内在属性必须在符号置换下保持不变。通过构造具体的符号置换和语义重赋值，不可公度性在一种解释下存在，在另一种解释下消失，故不可能是内在属性。这一定理构成了整个维度退化理论的元理论基石。它直接解释了为什么经典NP定义将非确定性的本质外包给外部量词``存在见证''和外部谓词$IsWitness(c)$，并在此过程中导致了不可公度性的丢失和对偶性的消除——因为经典确定性图灵机的内部语义无法表达刻画非确定性分支独立性（不可公度性）和结构（对偶性）的核心概念。
	
	\textbf{第二，复布尔图灵机（CBTM）。} 基于上述诊断，我们构建了CBTM，将计算符号从布尔值提升为四元伽罗瓦域$\mathbb{F}_4$中的代数元素。CBTM的核心洞见在于：非确定性计算对应于代数域的扩张——当读取虚部为$1$的符号时，计算必须分叉为两条路径，分别绑定于在基域$\mathbb{F}_2$上不可公度且互为对偶的生成元$\alpha$和$\beta$。通过投影算子$\operatorname{Re}$和$\operatorname{Im}$，我们分离``旧数据''与``新维度''，并引入双带视角将抽象的代数符号直观地分解为实带（确定性计算）和虚带（非确定性控制）。参数化等价定理（定理\ref{thm:equivalence}）严格确立了CBTM与经典模型的外延等价（$\mathbf{P}_{cb}=\mathbf{P}$，$\mathbf{NP}_{cb}=\mathbf{NP}$），同时揭示了CBTM操作语义的内涵增强——不可公度性内建、对偶性显式保留、拷贝操作被消除、语义内建而非外包——系统性地修正了经典NP定义的操作不自足和三重概念错误。参数化猜想将$\mathbf{P}$ vs.\ $\mathbf{NP}$等价转化为虚部$b$是否可消除的代数问题：域扩张$\mathbb{F}_4/\mathbb{F}_2$是否为非确定性计算提供了本质的代数资源？
	
	\textbf{第三，实布尔图灵机（RBTM）与生成元无关性。} 我们将抽象生成元$\alpha$具体化为代数数$\sqrt{2}$，引入RBTM和双带视角，将每个符号分解为实带（确定性数据）和虚带（非确定性控制）。虚带上的``1''直观地标记了``新维度''出现的位置，为后续的动态维度跟踪奠定了物理基础。生成元无关性定理（定理\ref{thm:generator-independence}）证明了计算能力不依赖于生成元的具体选择，揭示了非确定性的本质在于``引入与基域不可公度的新元素''这一事实本身，以及生成元之间的对偶关系。在此基础上，我们揭示了伽罗瓦自同构的操作性本质——路径切换——并据此提出敏感性猜想（猜想\ref{conj:sensitivity}）：P与NP的区分等价于是否存在某个非确定性选择点，使得单独翻转该点的选择能够改变计算结果。该猜想等价于$\mathbf{P}\neq\mathbf{NP}$（定理\ref{thm:sensitivity-equivalence}）。这一结论将P vs.\ NP问题从定量资源分析转化为定性结构判据，这一判据建立在对偶性的操作化之上——单点敏感性正是对偶性在单个选择点上的体现。
	
	\textbf{第四，方法论原则。} 我们提出了``数计一体''的方法论原则——计算理论应当仿照数学，要求任意算法都应当具有明确的语义，并且语义与语法保持一致，且操作语义必须是自足的。CBTM/RBTM框架通过参数化设计——不改变TM或NTM的代码，仅做一个巧妙的一一映射——将这一原则严格实现，系统性地修正了经典NP定义的操作不自足和三重概念错误，为计算理论从语法范式向语义范式的转变提供了模型基础。
	
	\subsection{从静态到动态：走向IVM与本质维度}
	
	本文所建立的CBTM/RBTM框架为非确定性计算提供了代数语义基础，但其代数结构是\textbf{静态}的。在CBTM中，无论计算过程中发生多少次非确定性分支，所有符号的虚部非零部分都只能关联于$\alpha$或$\beta$，整个计算所张成的代数空间维数最多为$2$。因此，CBTM虽然完美地刻画了单次分支的独立性（不可公度性）和对偶性，但无法区分``发生一次分支''与``发生了$n$次分支''——多次分支的代数维度会坍缩回一个常数（第5.8节）。RBTM同样面临这一局限：所有分支共享同一个生成元$\sqrt{2}$，生成元提取算子无法区分不同分支步骤所引入的维度。
	
	这一局限正是引入动态的虚部验证机（IVM）的根本动机。在后续论文中，我们将构建IVM作为CBTM的语义展开器。IVM的核心机制是动态生成元管理：借助生成元无关性定理所赋予的自由，为每次非确定性分支分配一个独立的素数平方根生成元$\sqrt{p_i}$，同时对偶地引入$-\sqrt{p_i}$。由素数平方根线性无关定理，这些生成元彼此线性无关，从而使得代数维度能够随分支次数线性累积，彻底解决静态框架的维度坍缩问题。基于IVM，我们将定义本质维度$\kappa(L)$——验证语言$L$所需的最小最坏情况代数维数——并证明它是一个障碍无关的不变量，能够区分$\mathbf{P}$与$\mathbf{NP}$而不受三大障碍干扰。
	
	结合本文提出的敏感性猜想，本质维度的下界证明将直接导向$\mathbf{P}\neq\mathbf{NP}$的证明：NP完全语言中不可消除的本质非确定性选择点，恰好对应于本质维度的不可消除性——需要动态引入的生成元数量无法被多项式界限压缩，从而$\kappa(L) = \Omega(n)$。而$\mathbf{P}$类语言可由零维验证器判定（$\kappa=0$），因此$\mathbf{P} \neq \mathbf{NP}$。
	
	这一从诊断（不可公度性元定理、操作不自足与三重概念错误）到模型构建（CBTM/RBTM，修正操作不自足和概念错误），再到度量工具（IVM和本质维度），最终到分离证明（$\mathbf{P}\neq\mathbf{NP}$）的完整逻辑链条，构成了维度退化理论的核心架构。本文为该架构奠定了模型基础和方法论原则。
	
	\appendix
	
	\section{NTM二元标准化的详细构造}
	
	\begin{proof}[引理\ref{lem:binary-ntm}的详细构造]
		设$N = (Q, \Sigma, \Gamma, \Delta, q_0, q_{\text{accept}}, q_{\text{reject}})$为一台多项式时间NTM，其转移关系在每一步至多产生$k \ge 2$个后继选择。构造$N' = (Q', \Sigma, \Gamma, \Delta', q_0, q_{\text{accept}}, q_{\text{reject}})$如下。
		
		对于$N$中的每个状态$q \in Q$和符号$a \in \Gamma$，设$\Delta(q,a) = \{(q_1,a_1,d_1), \ldots, (q_k,a_k,d_k)\}$。引入新的中间状态$q^{(1)}, \ldots, q^{(k-2)}$（均为新状态，不出现在$Q$中），并将$N$的一步$k$-way分支替换为以下结构：
		
		\begin{itemize}
			\item 在状态$q$读取$a$时，产生两个后继：$(q_1,a_1,d_1)$和$(q^{(1)},a,\text{stay})$。
			\item 在状态$q^{(1)}$读取$a$时，产生两个后继：$(q_2,a_2,d_2)$和$(q^{(2)},a,\text{stay})$。
			\item $\cdots$
			\item 在状态$q^{(k-2)}$读取$a$时，产生两个后继：$(q_{k-1},a_{k-1},d_{k-1})$和$(q_k,a_k,d_k)$。
		\end{itemize}
		
		在每个中间状态$q^{(j)}$，带头保持不动（stay操作可通过先右移再左移实现，或直接允许stay作为第三种带头移动方向）。这样，经过至多$k-1$步（每步产生两个分支），所有$k$个原选择被隔离到唯一的叶节点。由于$k$是常数，引入的状态数为$O(k)$，模拟时间的膨胀因子为$O(k) = O(1)$。因此$N'$满足引理的所有条件。对于不存在后继或仅有一个后继的格局，$N'$的行为与$N$相同（确定性转移或停机）。
	\end{proof}
	
	\section{参数化等价定理的详细证明}
	
	\begin{proof}[定理\ref{thm:equivalence}的详细证明]
		\textbf{(1) $\text{CBTM}|_0 \equiv_{\text{poly}} \text{DTM}$。}
		
		($\Rightarrow$) 设$M$为$\text{CBTM}|_0$。构造DTM $M'$如下：$M'$的状态集为$Q \times \{0,1\}$，其中第二个分量记录当前读取符号的虚部（始终为$0$）。纸带字母表为$\{0,1\}$。$M'$在每一步模拟$M$的确定性转移（$\Delta$中仅有一个匹配五元组），将CBTM的符号$\tau$映射为$\operatorname{Re}(\tau) \in \{0,1\}$，执行相同的带头移动。由于所有虚部为$0$，写入符号的实部由投影约束公理确定为布尔值，直接写入即可。模拟每步开销为$O(1)$，总时间为$O(T(n))$。
		
		($\Leftarrow$) 设$M$为经典DTM。构造$\text{CBTM}|_0$ $M'$如下：$M'$的状态集为$Q$。对于$M$的每个转移$\delta(q,a) = (q',a',d)$（其中$a,a' \in \{0,1\}$），$M'$的转移关系$\Delta$中包含五元组$(q, a+0\cdot\alpha, q', a'+0\cdot\alpha, d)$。由于所有符号虚部为$0$，分支触发公理要求$\Delta$中每个$(q,\tau)$仅有一个匹配五元组，满足条件。$M'$的纸带符号为$\mathbb{F}_2$中元素的自然嵌入（虚部为$0$）。模拟每步开销为$O(1)$，总时间为$O(T(n))$。
		
		\textbf{(2) $\text{CBTM} \equiv_{\text{poly}} \text{NTM}$。}
		
		($\Rightarrow$) 设$M$为CBTM。构造NTM $N$如下：$N$的状态集为$Q$，纸带字母表为$\{0,1\}$。$N$的转移关系$\Delta_N$包含：对于$M$的每个五元组$(q,\tau,q',\tau',d) \in \Delta$，在$\Delta_N$中添加$(q,\operatorname{Re}(\tau),q',\operatorname{Re}(\tau'),d)$。当$M$在读取符号$\tau$时触发分支（$\operatorname{Im}(\tau)=1$），$\Delta$中有两个五元组，对应地$\Delta_N$中在$(q,\operatorname{Re}(\tau))$下有两个五元组，$N$产生两个后继。模拟每步开销为$O(1)$，总时间为$O(T(n))$。
		
		($\Leftarrow$) 设$N$为经典NTM。由引理\ref{lem:binary-ntm}，可假设$N$在每个非确定性步骤恰好有两个选择。构造CBTM $M$如下：$M$的状态集为$Q$，纸带字母表包含$\{0,1,\alpha,\beta\}$。对于$N$的每个确定性转移（唯一五元组）$(q,a,q',a',d)$，在$\Delta$中添加$(q, a+0\cdot\alpha, q', a'+0\cdot\alpha, d)$。对于$N$的每个非确定性转移（两个五元组）$(q,a,q_1,a_1,d_1)$和$(q,a,q_2,a_2,d_2)$，在$\Delta$中添加$(q, a+1\cdot\alpha, q_1, a_1+1\cdot\alpha, d_1)$和$(q, a+1\cdot\beta, q_2, a_2+1\cdot\beta, d_2)$。
		
		需要验证这一构造满足分支触发公理和投影约束公理。在确定性转移中，读取符号虚部为$0$，$\Delta$中恰好有一个匹配五元组。在非确定性转移中，读取符号虚部为$1$，$\Delta$中恰好有两个匹配五元组（分别绑定$\alpha$和$\beta$）。写入符号的虚部由转移规则显式指定（分别为$1$和$1$），且满足投影约束：实部为$a_i$，虚部为$1$，均由当前状态和读取符号的投影通过布尔函数决定。模拟每步开销为$O(1)$，总时间为$O(T(n))$。
	\end{proof}
	
	\section{生成元无关性定理的补充细节}
	
	\begin{proof}[定理\ref{thm:generator-independence}的补充细节]
		本附录补充状态转移对应的严格形式化。
		
		设$M_\gamma = (Q_\gamma, \Sigma_\gamma, \Gamma_\gamma, \delta_\gamma, q_{0\gamma}, F_\gamma, \epsilon)$和$M_\delta = (Q_\delta, \Sigma_\delta, \Gamma_\delta, \delta_\delta, q_{0\delta}, F_\delta, \epsilon)$分别为基于生成元$\gamma$和$\delta$的RBTM。定义$\iota: Q_\gamma \to Q_\delta$为恒等映射（状态集可视为相同），$\psi: \Gamma_\gamma \to \Gamma_\delta$为$\psi(a+b\gamma) = a+b\delta$。
		
		对于任意$q \in Q_\gamma$和$\tau = a+b\gamma \in \Gamma_\gamma$，令$\delta_\gamma(q,\tau) = \{(q_1,\gamma_1,d_1), \ldots, (q_r,\gamma_r,d_r)\}$，其中$r \in \{1,2\}$由$b$的值决定（$b=0$时$r=1$，$b=1$时$r=2$）。由投影约束公理，每个写入符号$\gamma_i$可写为$\gamma_i = f_{\mathfrak{Re}}(q,a) + \gamma \cdot f_{\mathfrak{Im}}(q,b)$，其中$f_{\mathfrak{Re}}, f_{\mathfrak{Im}}$是取值于$\{0,1\}$的布尔函数。
		
		在$M_\delta$中，读取符号$\psi(\tau) = a+b\delta$，其投影值为$\mathfrak{Re}_\delta(\psi(\tau)) = a$，$\mathfrak{Im}_\delta(\psi(\tau)) = b$。由投影约束公理，$\delta_\delta(q,\psi(\tau))$中的写入符号$\delta_i$满足$\delta_i = f_{\mathfrak{Re}}(q,a) + \delta \cdot f_{\mathfrak{Im}}(q,b)$。因此$\psi(\gamma_i) = \delta_i$，且$(q_i, \delta_i, d_i) \in \delta_\delta(q, \psi(\tau))$。转移关系的对应性得证。
		
		同构保持计算树的完整结构，包括接受/拒绝状态和运行时间。具体而言，对于$M_\gamma$上的任意计算路径，其在$(\iota,\psi)$下的像恰为$M_\delta$上的一条计算路径，且接受状态被保持（因为$\iota(F_\gamma)=F_\delta$）。反之亦然。由于模拟是逐步的且每步开销为$O(1)$，运行时间在多项式意义上等价。故$M_\gamma$与$M_\delta$识别相同的语言且在多项式时间意义上等价。
	\end{proof}
	
	\bibliographystyle{plain}
	\bibliography{../cb}
	
\end{document}